% Chapter 0: Foundations and Prerequisites
% A comprehensive systems-engineering survey uniting classical multi-agent theory,
% contemporary coding agent architectures, transformer execution physics, and the operating kernel gap.

\makeatletter
\expandafter\def\csname pdchapternumberofprereq\endcsname{0}
\expandafter\def\csname pdchaptertitleofprereq\endcsname{Foundations and Prerequisites}
\expandafter\def\csname pdchaptercolorofprereq\endcsname{pdslate}
\expandafter\def\csname pdchapterpartlabelofprereq\endcsname{Prerequisites \textperiodcentered\ Foundations}
\expandafter\def\csname pdchapterpartindexofprereq\endcsname{0}
\expandafter\def\csname pdchapterepigraphofprereq\endcsname{I conclude that there are two ways of constructing a software design: One way is to make it so simple that there are obviously no deficiencies and the other way is to make it so complicated that there are no obvious deficiencies. The first method is far more difficult.}
\expandafter\def\csname pdchapterepigraphsourceofprereq\endcsname{C.\,A.\,R. Hoare, \emph{The Emperor's Old Clothes} (1980 Turing Award Lecture)}
\expandafter\def\csname pdchapterquestionofprereq\endcsname{What physical and theoretical ground must an agent operating kernel stand on?}
\expandafter\def\csname pdchapteronelineofprereq\endcsname{Classical multi-agent theory, contemporary coding agent harnesses, transformer KV-cache limits, and the kernel gap.}

\providecommand{\pdchapteropeningprereq}{%
  \pdbookrail{Theoretical Foundations}{Before constructing an operating kernel for autonomous systems,
    we must establish the physical and theoretical ground upon which modern AI agents stand:
    classical intentional agency, the actor model, transformer KV-cache memory scaling, the closed-loop harness,
    and which guarantees require mediated effects and retained evidence.}%
}

\providecommand{\pdchapterhandoffprereq}{%
  \pdbookrail{What \pdchapref{swk}{The Single-Writer Kernel} receives}{%
    The theoretical foundations of intentional agency, protocol mechanisms, and the physical
    limits of autoregressive inference identify the assumptions needed to mediate
    conflicting effects and recover obligations after a crash. Part I now builds the physical substrate: the minimal deterministic
    single-writer SQLite/WAL kernel that serializes local mutation and survives process crashes.}%
}

% Citation registrations for Chapter 0
\ifdefined\pdbookreference
  \pdbookreference{agha1986actors}{282}{G.~Agha. \newblock \emph{Actors: A Model of Concurrent Computation in Distributed Systems}. \newblock MIT Press, Cambridge, MA, 1986.}{Agha 1986}
  \pdbookreference{akerlof1970}{283}{George A. Akerlof. \newblock The Market for ``Lemons'': Quality Uncertainty and the Market Mechanism. \newblock \textit{Quarterly Journal of Economics}, 84(3):488--500, 1970.}{Akerlof 1970}
  \pdbookreference{arrow1951social}{284}{K.~J. Arrow. \newblock \emph{Social Choice and Individual Values}. \newblock John Wiley \& Sons, New York, 1951.}{Arrow 1951}
  \pdbookreference{austin1962how}{285}{J.~L. Austin. \newblock \emph{How to Do Things with Words}. \newblock Oxford University Press, 1962.}{Austin 1962}
  \pdbookreference{bainbridge1983}{286}{Lisanne Bainbridge. \newblock Ironies of Automation. \newblock \textit{Automatica}, 19(6):775--779, 1983.}{Bainbridge 1983}
  \pdbookreference{becker1968}{287}{G.~S. Becker. \newblock Crime and punishment: An economic approach. \newblock \emph{Journal of Political Economy}, 76(2):169--217, 1968.}{Becker 1968}
  \pdbookreference{birgisson2014macaroons}{288}{Arnar Birgisson, Joe Gibbs Politz, Úlfar Erlingsson, Ankur Taly, Michael Vrable, and Mark Lentczner. \newblock Macaroons: Cookies with Contextual Caveats for Decentralized Authorization in the Cloud. \newblock In \textit{Network and Distributed System Security Symposium (NDSS)}, 2014. \newblock \href{https://doi.org/10.14722/ndss.2014.23212}{Online source}}{Birgisson et al. 2014}
  \pdbookreference{bordini2007agentspeak}{289}{R.~H. Bordini, J.~F. H{\"u}bner, and M.~Wooldridge. \newblock \emph{Programming Multi-Agent Systems in AgentSpeak using Jason}. \newblock John Wiley \& Sons, 2007.}{Bordini et al. 2007}
  \pdbookreference{bratman1987intention}{290}{M.~E. Bratman. \newblock \emph{Intention, Plans, and Practical Reason}. \newblock Harvard University Press, Cambridge, MA, 1987.}{Bratman 1987}
  \pdbookreference{charrier2014bigbrother}{292}{T.~Charrier, F.~Ouchet, and F.~Schwarzentruber. \newblock \emph{Big Brother Logic: Reasoning about Agents Equipped with Surveillance Cameras in the Plane (Demonstration)}. \newblock In \emph{Proc. AAMAS}, pages 1633--1634, 2014. \newblock \href{https://aamas.csc.liv.ac.uk/Proceedings/aamas2014/aamas/p1633.pdf}{Primary record}.}{Charrier et al. 2014}
  \pdbookreference{clarke1971multipart}{293}{E.~H. Clarke. \newblock \emph{Multipart pricing of public goods}. \newblock \emph{Public Choice}, 11(1):17--33, 1971.}{Clarke 1971}
  \pdbookreference{fipa2002acl}{294}{FIPA. \newblock \emph{FIPA ACL Communicative Act Library Specification}. \newblock Technical Report SC00037J, Foundation for Intelligent Physical Agents, 2002.}{FIPA 2002}
  \pdbookreference{gibbard1973manipulation}{295}{A.~Gibbard. \newblock \emph{Manipulation of voting schemes: A general result}. \newblock \emph{Econometrica}, 41(4):587--601, 1973.}{Gibbard 1973}
  \pdbookreference{groves1973incentives}{296}{T.~Groves. \newblock \emph{Incentives in teams}. \newblock \emph{Econometrica}, 41(4):617--631, 1973.}{Groves 1973}
  \pdbookreference{hewitt1973actor}{297}{C.~Hewitt, P.~Bishop, and R.~Steiger. \newblock \emph{A universal modular ACTOR formalism for artificial intelligence}. \newblock In \emph{Proc. 3rd International Joint Conference on Artificial Intelligence (IJCAI)}, pages 235--245, 1973.}{Hewitt et al. 1973}
  \pdbookreference{hoare1978csp}{298}{C.~A.~R. Hoare. \newblock \emph{Communicating sequential processes}. \newblock \emph{Communications of the ACM}, 21(8):666--677, 1978.}{Hoare 1978}
  \pdbookreference{lamport1978}{299}{Leslie Lamport. \newblock Time, Clocks, and the Ordering of Events in a Distributed System. \newblock \textit{Communications of the ACM}, 21(7):558--565, 1978.}{Lamport 1978}
  \pdbookreference{liu2024lost}{300}{Nelson F.\ Liu, Kevin Lin, John Hewitt, Ashwin Paranjape, Michele Bevilacqua, Fabio Petroni, and Percy Liang. \newblock Lost in the Middle: How Language Models Use Long Contexts. \newblock \textit{TACL}, 12:157--173, 2024. arXiv:2307.03172.}{Liu et al. 2024}
  \pdbookreference{myerson1983efficient}{301}{R.~B. Myerson and M.~A. Satterthwaite. \newblock \emph{Efficient mechanisms for bilateral trading}. \newblock \emph{Journal of Economic Theory}, 29(2):265--281, 1983.}{Myerson \& Satterthwaite 1983}
  \pdbookreference{ostrom1990governing}{302}{Elinor Ostrom. \newblock \textit{Governing the Commons: The Evolution of Institutions for Collective Action}. \newblock Cambridge University Press, 1990.}{Ostrom 1990}
  \pdbookreference{ramadge1987supervisory}{303}{P.~J.~G. Ramadge and W.~M. Wonham. \newblock \emph{Supervisory control of a class of discrete event processes}. \newblock \emph{SIAM Journal on Control and Optimization}, 25(1):206--230, 1987.}{Ramadge \& Wonham 1987}
  \pdbookreference{rao1991modeling}{304}{Anand~S. Rao and Michael~P. Georgeff. \newblock Modeling Rational Agents within a BDI-Architecture. \newblock In \textit{International Conference on Principles of Knowledge Representation and Reasoning (KR)}, pages 473--484, 1991.}{Rao \& Georgeff 1991}
  \pdbookreference{rao1995bdi}{305}{A.~S. Rao and M.~P. Georgeff. \newblock \emph{BDI agents: From theory to practice}. \newblock In \emph{Proc. 1st International Conference on Multi-Agent Systems (ICMAS)}, pages 312--319, 1995.}{Rao \& Georgeff 1995}
  \pdbookreference{saltzer1975protection}{306}{Jerome H. Saltzer and Michael D. Schroeder. \newblock The Protection of Information in Computer Systems. \newblock \textit{Proceedings of the IEEE}, 63(9):1278--1308, 1975.}{Saltzer \& Schroeder 1975}
  \pdbookreference{satterthwaite1975strategy}{307}{M.~A. Satterthwaite. \newblock \emph{Strategy-proofness and Arrow's conditions: Existence and correspondence theorems for voting procedures and social welfare functions}. \newblock \emph{Journal of Economic Theory}, 10(2):187--217, 1975.}{Satterthwaite 1975}
  \pdbookreference{searle1969speech}{308}{J.~R. Searle. \newblock \emph{Speech Acts: An Essay in the Philosophy of Language}. \newblock Cambridge University Press, 1969.}{Searle 1969}
  \pdbookreference{shinn2023reflexion}{309}{N.~Shinn, F.~Cassano, E.~Berman, A.~Gopinath, K.~Narasimhan, and S.~Yao. \newblock \emph{Reflexion: Language Agents with Verbal Reinforcement Learning}. \newblock In \emph{Proc. NeurIPS}, 2023. \newblock \href{https://arxiv.org/abs/2303.11366v4}{Primary record}.}{Shinn et al. 2023}
  \pdbookreference{shoham2009multiagent}{310}{Y.~Shoham and K.~Leyton-Brown. \newblock \emph{Multiagent Systems: Algorithmic, Game-Theoretic, and Logical Foundations}. \newblock Cambridge University Press, 2009.}{Shoham \& Leyton-Brown 2009}
  \pdbookreference{smith1980cnp}{311}{R.~G. Smith. \newblock \emph{The Contract Net Protocol: High-level communication and control in a distributed problem solver}. \newblock \emph{IEEE Transactions on Computers}, C-29(12):1104--1113, 1980.}{Smith 1980}
  \pdbookreference{vickrey1961counterspeculation}{312}{W.~Vickrey. \newblock \emph{Counterspeculation, auctions, and competitive sealed tenders}. \newblock \emph{The Journal of Finance}, 16(1):8--37, 1961.}{Vickrey 1961}
  \pdbookreference{yao2022react}{313}{Yao et al. 2023. \emph{ReAct: Synergizing reasoning and acting in language models}.}{Yao et al. 2023}
  \pdbookreference{yao2023tree}{314}{Yao et al. 2023. \emph{Tree of Thoughts: Deliberate problem solving with large language models}.}{Yao et al. 2023}
  \pdbookreference{gui2025hypertree}{315}{R.~Gui, Z.~Wang, J.~Wang, C.~Ma, H.~Zhen, M.~Yuan, J.~Hao, D.~Lian, E.~Chen, and F.~Wu. \newblock \emph{HyperTree Planning: Enhancing LLM Reasoning via Hierarchical Thinking}. \newblock arXiv:2505.02322v2, 2025. \newblock \href{https://arxiv.org/abs/2505.02322v2}{Primary record}.}{Gui et al. 2025}
  \pdbookreference{conway1968}{316}{M.~E. Conway. \newblock \emph{How do Committees Invent?} \newblock \emph{Datamation}, 14(4):28--31, 1968.}{Conway 1968}
  \pdbookreference{dao2022flashattention}{317}{T.~Dao, D.~Y. Fu, S.~Ermon, A.~Rudra, and C.~R{\'e}. \newblock \emph{FlashAttention: Fast and Memory-Efficient Exact Attention with IO-Awareness}. \newblock In \emph{Proc. 36th Conference on Neural Information Processing Systems (NeurIPS)}, 2022.}{Dao et al. 2022}
  \pdbookreference{hong2024metagpt}{318}{S.~Hong, M.~Zhuge, J.~Chen, X.~Zheng, Y.~Cheng, C.~Zhang, J.~Wang, Z.~Wang, S.~K.~S. Yau, Z.~Lin, et al. \newblock \emph{MetaGPT: Meta Programming for Multi-Agent Collaborative Framework}. \newblock In \emph{Proc. International Conference on Learning Representations (ICLR)}, 2024.}{Hong et al. 2024}
  \pdbookreference{patel2026sixsigma}{319}{K.~Patel, S.~Surendira, J.~George, and S.~Kapale. \newblock \emph{The Six Sigma Agent: Achieving Enterprise-Grade Reliability in LLM Systems Through Consensus-Driven Decomposed Execution}. \newblock arXiv:2601.22290v1, 2026. \newblock \href{https://arxiv.org/abs/2601.22290v1}{Primary record}.}{Patel et al. 2026}
  \pdbookreference{kwon2023vllm}{320}{W.~Kwon, Z.~Li, S.~Zhuang, Y.~Sheng, L.~Zheng, C.~H. Yu, J.~E. Gonzalez, H.~Zhang, and I.~Stoica. \newblock \emph{Efficient Memory Management for Large Language Model Serving with PagedAttention}. \newblock In \emph{Proc. 29th ACM Symposium on Operating Systems Principles (SOSP)}, pages 611--626, 2023.}{Kwon et al. 2023}
  \pdbookreference{liang2022helm}{321}{P.~Liang, R.~Bommasani, T.~Lee, D.~Tsipras, D.~Soylu, M.~Yasunaga, Y.~Zhang, D.~Narayanan, Y.~Wu, A.~Kumar, et al. \newblock \emph{Holistic Evaluation of Language Models}. \newblock \emph{Transactions on Machine Learning Research (TMLR)}, 2023.}{Liang et al. 2023}
  \pdbookreference{liu2023agentbench}{322}{X.~Liu, H.~Yu, H.~Zhang, Y.~Xu, X.~Lei, H.~Lai, Y.~Gu, H.~Ding, K.~Men, K.~Yang, et al. \newblock \emph{AgentBench: Evaluating LLMs as Agents}. \newblock In \emph{Proc. International Conference on Learning Representations (ICLR)}, 2024.}{Liu et al. 2024}
  \pdbookreference{ouyang2022instruct}{323}{L.~Ouyang, J.~Wu, X.~Jiang, D.~Almeida, C.~L. Wainwright, P.~Mishkin, C.~Zhang, S.~Agarwal, K.~Slama, A.~Ray, et al. \newblock \emph{Training language models to follow instructions with human feedback}. \newblock In \emph{Proc. 36th Conference on Neural Information Processing Systems (NeurIPS)}, 2022.}{Ouyang et al. 2022}
  \pdbookreference{rafailov2023dpo}{324}{R.~Rafailov, A.~Sharma, E.~Mitchell, S.~Ermon, C.~D. Manning, and C.~Finn. \newblock \emph{Direct Preference Optimization: Your Language Model is Secretly a Reward Model}. \newblock In \emph{Proc. 37th Conference on Neural Information Processing Systems (NeurIPS)}, 2023.}{Rafailov et al. 2023}
  \pdbookreference{vaswani2017attention}{325}{A.~Vaswani, N.~Shazeer, N.~Parmar, J.~Uszkoreit, L.~Jones, A.~N. Gomez, L.~Kaiser, and I.~Polosukhin. \newblock \emph{Attention Is All You Need}. \newblock In \emph{Proc. 31st Conference on Neural Information Processing Systems (NeurIPS)}, pages 5998--6008, 2017.}{Vaswani et al. 2017}
  \pdbookreference{wang2023voyager}{326}{Wang et al. 2023. \emph{Voyager: An Open-Ended Embodied Agent with Large Language Models}.}{Wang et al. 2023}
  \pdbookreference{wang2024tdag}{327}{Y.~Wang, Z.~Wu, J.~Yao, and J.~Su. \newblock \emph{TDAG: A Multi-Agent Framework based on Dynamic Task Decomposition and Agent Generation}. \newblock arXiv:2402.10178, 2024; revised 2025. \newblock \href{https://arxiv.org/abs/2402.10178v2}{Primary record}.}{Wang et al. 2024}
  \pdbookreference{wang2023agentsurvey}{328}{L.~Wang, C.~Ma, X.~Feng, Z.~Zhang, H.~Yang, J.~Zhang, Z.~Chen, J.~Tang, X.~Chen, Y.~Lin, W.~X. Zhao, Z.~Wei, and J.-R. Wen. \newblock \emph{A Survey on Large Language Model based Autonomous Agents}. \newblock arXiv:2308.11432, 2023; revised 2025. \newblock \href{https://arxiv.org/abs/2308.11432v7}{Primary record}.}{Wang et al. 2023}
  \pdbookreference{zhou2026formaljudge}{329}{J.~Zhou, Y.~Sheng, H.~Lou, Y.~Yang, and J.~Fu. \newblock \emph{FormalJudge: A Neuro-Symbolic Paradigm for Agentic Oversight}. \newblock arXiv:2602.11136v2, 2026. \newblock \href{https://arxiv.org/abs/2602.11136v2}{Primary record}.}{Zhou et al. 2026}
  \pdbookreference{anthropic2026hooks}{330}{Anthropic. \newblock \emph{Hooks reference}. \newblock Claude Code documentation, accessed 3 November 2026. \newblock \href{https://code.claude.com/docs/en/hooks}{Official documentation}.}{Anthropic 2026}
  \pdbookreference{jasonfaq}{331}{Jason developers. \newblock \emph{Jason FAQ: AgentSpeak Language}. \newblock Official documentation, accessed 26 October 2026. \newblock \href{https://jason-lang.github.io/doc/faq.html}{Primary documentation}.}{Jason developers, accessed 2026}
  \pdbookreference{jasonfailureimplementation}{332}{Jason developers. \newblock \emph{Jason interpreter: goal-failure handling}. \newblock Source revision \texttt{4b98ddb1f9c8}, accessed 26 October 2026. \newblock \href{https://github.com/jason-lang/jason/blob/4b98ddb1f9c88d3ecdd84c6559f8effbf20f8287/jason-interpreter/src/main/java/jason/asSemantics/TransitionSystem.java}{Pinned source}.}{Jason developers, accessed 2026}
\fi
\makeatother
\pdchapter{0}{\texorpdfstring{Foundations and\\[2pt]Prerequisites}{Foundations and Prerequisites}}{prereq}{pdslate}
% The display title breaks on the opener; the running head stays on one line.
\fancyhead[RE,LO]{\small\itshape\color{hhgray} Chapter 0. Foundations and Prerequisites}

\pdchapteropeningprereq

\ifdefined\pdswisslandmarkplate
\begin{figure}[p]
  \centering
  \pdswisslandmarkplate{01}{\linewidth}
\end{figure}
\fi

\setlength{\emergencystretch}{2em}

\noindent\textbf{Start with one interrupted job.}
Alice asks two workers to fix the same checkout failure. Both read revision~7.
The first publishes a repair; the second overwrites it with a patch based on
revision~7. Then both processes exit. Alice needs three answers: which change
was permitted, which change took effect, and which work is still owed?

\begin{pdexample}{Two increments, one missing update}
Start a counter at 0. Each worker reads 0, computes $0+1$, and writes 1.
The final counter is 1. If each entire read--increment--write operation runs
before the other begins, the states are $0\to1\to2$. Merely putting the two
final writes in a queue still gives $0\to1\to1$. We will return to this
stale-observation problem when we define admission and recovery.
\end{pdexample}

\noindent\textbf{The route.}
First distinguish a worker's proposal, an admitted effect and a retained
record. Then examine commitments, messages and incentives; account for the
cost of keeping information in a prompt; and follow a request through a
refusal, a timeout and a restart. The final result assembles these cases.
Read the examples before the formalism. The literature comparisons identify
obligations that return in the later chapters, not a vocabulary to memorize.

% Flush the opening art before the first definition, never through its paragraph.
\ifdefined\pdswisslandmarkplate\clearpage\fi
\section{Introduction: The Unit of Account, Agency, and the Kernel Gap}

\subsection{The Unit of Account is the Work}

A \emph{work unit} is the case file for a job: who requested it, what actions
are allowed, what has been observed, and what remains to be done. An
\emph{obligation} is a recorded requirement, such as supplying a test result
before the job may be called complete. Alice's repair remains the same job
when a worker exits. Chapters 1 and 5 develop its record and its succession.

\subsection{Defining the Agent and the Kernel}
The counter example needs a proposer, an actuator and a rule separating them.
An \emph{actuator} actually changes a resource, for example by writing a file.
A proposal may describe a write without yet performing it.

\begin{definition}[Autonomous Agent]\label{def:agent}
An \textbf{agent} chooses a proposed action $a_t\in\mathcal A$ after an
observation history $h_t\in\mathcal H$. Its policy
$\pi_\theta(a_t\mid h_t)$ specifies the probability of that choice;
$\theta$ denotes its parameters or rules. A deterministic policy assigns
probability one to its chosen action. The surrounding harness runs the
policy, invokes tools and returns observations.
\end{definition}
For Alice's worker, a patch is a proposal and a test result is an observation.
A passive log that records requests but makes no choices is not an agent in
this model. Neither a proposal nor a reported belief certifies the host state.

\begin{definition}[Operating Kernel]\label{def:kernel}
The \textbf{coordination kernel} is the trusted application component that
admits changes to specified resources and retains coordination state. Its
contract names the operations it mediates, the invariants it checks, the
evidence it retains and the failures its storage survives. A state invariant
is a property required before and after each admitted transition, such as
at most one current holder of a particular lease.
\end{definition}
Here ``kernel'' does not mean the operating-system kernel. A daemon (a
background service) and its database implement Chapter 1's construction.
Host isolation and external actuators need additional controls. A log written
after an unrestricted file write is evidence, not an admission gate; a
coordination-database transaction does not make that file write atomic.

\subsection{The Dual Lineage of Autonomous Agency}\label{sec:prereq-lineage}
Symbolic multi-agent research supplies explicit representations of beliefs,
goals, plans and messages. Generative models supply another way to propose
plans and interpret unfamiliar observations. Neither representation settles
who may overwrite revision~7. That depends on the surrounding protocol.

\subsection{Proposals and checked effects}
\begin{table}[H]
\centering\small
\begin{tabularx}{\textwidth}{@{}lXX@{}}
\toprule
\textbf{Question} & \textbf{Proposal supplies} & \textbf{System checks}\\
\midrule
What should change? & A patch & Authorized resource and revision\\
What is still owed? & A remembered plan & Durable open obligations\\
Did it work? & A reported outcome & Specified evidence\\
May another worker act? & A delegation request & Current scope and expiry\\
\bottomrule
\end{tabularx}
\caption{The distinction applies to symbolic and neural proposers alike.}
\label{tab:lineage-comparison}
\end{table}
The last section proves an elementary preservation claim for the checked
boundary. Chapter 2 returns to the last row: communicating a task description
and granting permission to act are different operations.

\subsection{Eight Propositions of Accountable Coordination}

These previews introduce the questions that recur in the eight chapters. They are not eight proofs or claims that the complete system exists. Each owning chapter supplies its definitions, examples and limits.

\subsubsection*{Proposition 1: Prevention Begins at the Effect Boundary (\pdchapref{swk}{The Single-Writer Kernel})}
Preventing a forbidden effect requires control over every path that can produce it. Advisory instructions alone do not provide that control. For example, if two agents can read the same value $x=0$ and later write their independently computed value $x+1$, the schedule read--read--write--write ends at $x=1$, losing one increment. This is a permitted counterexample, not an inevitable outcome of concurrency: disjoint updates and correctly implemented transactions can be safe.

The Single-Writer Kernel checks a request's authorization and state preconditions at the same boundary that commits its ledger transition. A complete read--modify--write transaction, or an atomic version check followed by the write, rejects the stale update. Serializing only the final writes is insufficient. SQLite's \href{https://www.sqlite.org/isolation.html}{transaction isolation} covers its database operations; making a filesystem or network effect atomic with a receipt requires a separate actuator and crash-recovery protocol. Strict serializability additionally requires that the operation order respect real-time completion.

\begin{figure}[tp]
  \centering
  % Reader question: Where can a stale write be refused?
% Claim: A transaction or version guard covers the full read-modify-write operation.
% Evidence: Two clients read version7; only the first version-checked mutation commits.
% Distinguish: Proposal, admission and durable commit.
% Grammar: A sequence, not a funnel suggesting WAL alone prevents lost updates.
% Acceptance: Identify why client B is refused.
\pdfigurehue{pdcobalt}%
\begin{tikzpicture}[pd figure,x=1cm,y=1cm]

  \node[pd artifact,text width=3.4cm,align=center] (a) at (2.4,0) {Client A\\Read version 7; propose 8};
  \node[pd artifact,text width=3.4cm,align=center] (b) at (8.0,0) {Client B\\Read version 7; propose 8};
  \node[pd focus state,text width=6.4cm,align=center] (gate) at (5.2,-1.8) {Commit boundary\\Atomically check expected version and write};
  \node[pd ready state,text width=3.4cm,align=center] (commit) at (2.4,-3.6) {A admitted\\Version 8 durably committed};
  \node[pd breach state,text width=3.4cm,align=center] (reject) at (8.0,-3.6) {B refused\\Expected 7; actual 8};
  \draw[pd arrow] (a) -- (gate);
  \draw[pd arrow] (b) -- (gate);
  \draw[pd ready arrow] (gate) -- (commit);
  \draw[pd breach arrow] (gate) -- (reject);
  \node[pd note,text width=9.4cm,align=left] at (5.2,-5.3) {A serialized log records the accepted order. Preventing lost updates also requires a transaction or version check covering each complete operation.};

\end{tikzpicture}
%
  \pdwidecaption{\textbf{Serial admission with version validation.} Two proposals based on version 7 cannot both replace it. The gate commits one and rejects the stale request. The guarantee requires validation and effect to share one indivisible boundary.}[fig:spark-prop1]
\end{figure}

Formally, let $e_i$ and $e_j$ be concurrent events under Lamport's happened-before relation~\pdcite{lamport1978}. Their incomparability alone does not imply an error. A conflict arises when their resource accesses interfere and at least one writes. An admission gate assigns admitted operations a commit order $\sigma(e)$, checks each operation against the current version, and makes the check and state transition indivisible. Under complete mediation and correct recovery, this preserves any invariant maintained by each admitted transition. Locks, compare-and-swap and transactional validation are alternative ways to implement that boundary.

\subsubsection*{Proposition 2: Delegation Only Narrows (\pdchapref{anchor}{The Anchor Protocol})}
In classical security engineering, authority must obey Saltzer and Schroeder's fundamental Principle of Least Privilege~\pdcite{saltzer1975protection}: every process must operate using the least set of privileges necessary to complete its assigned task. In naive multi-agent architectures where agents delegate subtasks by synthesizing unconstrained prompts, authority frequently inflates due to broad, ambiguous system instructions, leading to catastrophic privilege escalation.

Under the Anchor Protocol, every child capability has authority and expiry no greater than its parent's ($C_0 \supseteq C_1 \supseteq C_2$ and $T_0 \ge T_1 \ge T_2$). Authentication proves who signed the credential; attenuation specifies what the signature is permitted to authorize. Chained cryptographic capability tokens (such as macaroons~\pdcite{birgisson2014macaroons} and Harbor Cards) enforce monotonic restriction at every delegation hop: an intermediate delegate cannot widen permissions, extend deadlines, or remove sandbox constraints when dispatching subordinate workers.

\begin{figure}[tp]
  \centering
  % Reader question: What can a delegate pass onward?
% Claim: Every accepted child grant is a subset of its parent grant.
% Evidence: An illustrative path scope and expiry chain.
% Distinguish: Permitted narrowing and rejected widening.
% Grammar: A delegation chain with a refused branch.
% Acceptance: Trace which child request violates attenuation.
\pdfigurehue{pdcobalt}%
\begin{tikzpicture}[pd figure,x=1cm,y=1cm]

  \node[pd identity state,text width=6.0cm,align=center] (parent) at (4.5,0) {Parent grant\\Read and write repository; expires 18:00};
  \node[pd focus state,text width=6.0cm,align=center] (child) at (4.5,-1.8) {Child grant\\Read and write docs; expires 17:00};
  \node[pd ready state,text width=6.0cm,align=center] (leaf) at (4.5,-3.6) {Accepted descendant\\Read docs only; expires 16:00};
  \draw[pd arrow] (parent) -- (child);
  \draw[pd ready arrow] (child) -- (leaf);
  \node[pd breach state,text width=6.0cm,align=center] (bad) at (4.5,-5.4) {Refused request\\Write source or extend expiry to 19:00};
  \draw[pd breach arrow] (child.east) -- (9.0,-1.8) -- (9.0,-5.4) -- (bad.east);

\end{tikzpicture}
%
  \pdwidecaption{\textbf{Capability attenuation.} Child scope and expiry do not exceed the parent's. Chain verification and conjunctive caveat checks enforce this relation; they do not revoke a broader ancestor token already held by a delegate.}[fig:spark-prop2]
\end{figure}

Mathematically, a capability token $M_k$ at delegation depth $k$ is constructed as an HMAC chain:
\begin{equation}
s_k \;=\; \mathrm{HMAC}(s_{k-1}, c_k)
\end{equation}
Here $s_0$ denotes the authenticated root tag derived with the issuer's secret, and $c_k$ is an added caveat. Under HMAC unforgeability, a holder of the final tag cannot manufacture a valid tag for a less restricted chain it was never issued. The verifier must authenticate the complete chain and conjoin every caveat. A holder that retained a broader ancestor token can still use that token until it expires or is revoked; attenuation does not revoke previously issued authority.

\subsubsection*{Proposition 3: Confinement Is a Property of the Room, Not of the Tenant (\pdchapref{sealed}{The Sealed Harbor})}
In multi-tenant, multi-principal systems, mutual distrust is the necessary default. Two owners who trust neither each other nor the hosting venue can still take a verified computation out of a sealed room, provided the execution chamber guarantees non-bypassable physical isolation. Confinement cannot be promised by the tenant model: an LLM process is vulnerable to prompt injection, memory corruption, or adversarial jailbreaking.

Confinement is supplied by a specified host policy and its trusted computing base. The intended room closes unauthorized network and filesystem paths and meters the admitted egress channel. Its guarantee covers the modeled channels only. Timing, termination, shared hardware and other side channels need separate analysis. Canary tokens provide evidence of particular monitored disclosures; they cannot establish that no other secret or encoding escaped.

\begin{figure}[tp]
  \centering
  % Reader question: Which path can an untrusted process use to leave its sandbox?
% Claim: A confinement claim depends on mediating all specified effect paths.
% Evidence: Sandbox, explicit egress boundary and configured checks.
% Distinguish: Mediated egress and excluded covert-channel claims.
% Grammar: Trust boundary diagram, not an absolute security shield.
% Acceptance: Name the enforced path and the remaining assumption.
\pdfigurehue{pdcobalt}%
\begin{tikzpicture}[pd figure,x=1cm,y=1cm]

  \node[pd state,text width=3.6cm,align=center] (agent) at (3.0,0) {Untrusted process\\Private scratch and memory};
  \node[pd trust boundary,fit=(agent),inner sep=12pt] (sandbox) {};
  \node[pd title,text width=4.0cm,align=center] at (3.0,1.65) {Sandbox boundary};
  \node[pd focus state,text width=4.0cm,align=center] (gate) at (3.0,-2.2) {Egress mediator\\Destination, budget and content checks};
  \node[pd protocol state,text width=4.0cm,align=center] (dest) at (3.0,-4.4) {Authorized destination};
  \node[pd breach state,text width=2.9cm,align=center] (deny) at (8.0,-2.2) {Refused\\No traffic leaves};
  \draw[pd arrow] (agent) -- (gate);
  \draw[pd ready arrow] (gate) -- (dest);
  \draw[pd breach arrow] (gate) -- (deny);
  \node[pd note,text width=9.2cm,align=left] at (5.0,-5.8) {Canary matching detects matching content. Timing channels, unmediated interfaces and arbitrary secrets require separate analysis.};

\end{tikzpicture}
%
  \pdwidecaption{\textbf{Mediated egress.} The room closes unauthorized modeled paths and admits traffic through a checked channel. Canary matching supplies evidence of particular disclosures. Timing and other unmodeled side channels require separate analysis.}[fig:spark-prop3]
\end{figure}

The confinement policy is expressed as a non-bypassable sandbox profile $P_{\text{room}}$ compiled into the kernel:
\begin{equation}
\forall \sigma \in \Sigma_{\text{syscall}}, \quad \mathrm{Permit}(\sigma \mid P_{\text{room}}) \iff \sigma \in \Sigma_{\text{confined}}
\end{equation}
Here $\Sigma_{\text{confined}}$ is shorthand for the allowed syscall-and-argument cases in a stated sandbox policy; executable identity, paths and inherited descriptors matter as well as syscall names. A canary $\kappa\sim\mathcal{U}(\{0,1\}^{256})$ may be registered with an egress matcher. Observing its registered representation is a disclosure witness under the matcher's attribution assumptions. An adversary that leaks a different secret or an unrecognized encoding may evade it. Revocation and any financial penalty therefore follow the declared evidence and contestation rules.

\subsubsection*{Proposition 4: Supervision Has an Information Cost (\pdchapref{ls}{Attention and Supervision})}
Human oversight is constrained by fundamental cognitive and information-theoretic limits. As Bainbridge observed in \emph{Ironies of Automation}~\pdcite{bainbridge1983}, automating routine operations does not eliminate human responsibility; it concentrates attention on rare, high-consequence failure modes. A human operator cannot parse a raw firehose of millions of agent telemetry tokens without suffering cognitive saturation and supervisory fatigue.

Yet naive LLM-generated summaries are actively dangerous: they hallucinate away subtle test failures, sanitize errors, and declare false victories. A digest cannot guarantee that a supervisor will notice any one of many relevant events unless it expends sufficient bits and attention to distinguish them. In an accountable architecture, a summary is therefore an index into evidence, not a substitute for it. The supervisory dashboard presents low-entropy summaries backed by cryptographic pointers directly into the underlying Merkle event log, allowing instant drilling down from high-level status to raw syscall traces.

\begin{figure}[tp]
  \centering
  % Reader question: how does a digest lead to retained execution evidence?
% Claim: four resolutions index the same evidence log; hashes do not recover bytes.
% Evidence: schematic references, with receipt R42 as the worked identifier.
% Distinguish: summary, commitment, retained artifacts and semantic correctness.
% Grammar: resolution stack plus shared log, native-size keyed lookup table.
% Acceptance: locate R42 and explain why its artifacts must still be retained.
\begin{minipage}{11.1cm}
\centering\pdfigurehue{pdcobalt}%
\begin{tikzpicture}[pd figure,x=1cm,y=1cm]
  \def\h{1.35}\def\w{6.10}\def\gx{-0.65}
  \def\lx{7.10}\def\lw{2.55}\def\stackH{5.40}
  \foreach \i/\name/\detail/\n in {
    3/{digest}/{one-line fleet summary}/1,
    2/{fleet status}/{records grouped by project}/2,
    1/{execution receipt}/{invocation R42 and its cost}/3,
    0/{syscall trace}/{retained open/read/write events}/4}{
    \draw[pd legible fill] (0,{\i*\h}) rectangle (\w,{(\i+1)*\h});
    \node[pd title,anchor=west] at (.20,{\i*\h+.98}) {\name};
    \node[pd label,anchor=west,text width=5.65cm,align=left]
      at (.20,{\i*\h+.40}) {\detail};
    \node[pd badge] at (\gx,{\i*\h+.67}) {\n};
    \draw[pd truth arrow] (\w,{\i*\h+.67})--(\lx,{\i*\h+.67});}
  \node[pd truth state,minimum width=\lw cm,minimum height=\stackH cm]
    at ({\lx+\lw/2},{\stackH/2}) {};
  \foreach \k in {1,...,7}{
    \draw[pd hairline] (\lx,{\k*\stackH/8})--({\lx+\lw},{\k*\stackH/8});}
  \node[pd title,anchor=south,text width=2.55cm,align=center]
    at ({\lx+\lw/2},{\stackH+.18}) {Merkle evidence log};
\end{tikzpicture}
\par\smallskip
{\pdfiglabel\setlength{\tabcolsep}{3pt}%
\begin{tabularx}{\linewidth}{@{}lX@{}}
\toprule
\textbf{Resolution} & \textbf{Retained support} \\\midrule
1 Digest & References to the evidence log as a whole. \\
2 Fleet status & The same records, grouped by project. \\
3 Receipt R42 & A committed leaf identifying one invocation. \\
4 Syscall trace & Trace bytes checked against the commitment. \\\bottomrule
\end{tabularx}}
\par\smallskip{\pdfiglabelitalic References navigate; commitments check retrieved bytes.
Neither reconstructs missing artifacts or proves their semantic correctness.\par}
\end{minipage}
%
  \pdwidecaption{\textbf{An index into retained evidence.} A status claim links to receipt R42 and its retained artifacts. The digest does not reconstruct omitted data; the evidence store must retain it.}[fig:spark-prop4]
\end{figure}

Let $X$ be a finite discrete execution history and $D=f(X)$ its digest. Then $I(X;D)=H(D)\le H(X)$. If $H(D)<H(X)$, the digest alone cannot reconstruct the history: $H(X\mid D)=H(X)-H(D)>0$. Retained receipts are the additional information that makes drill-down possible. A balanced binary Merkle tree over $N$ retained receipts supplies an inclusion path with $O(\log N)$ hashes. That bound concerns proof size and verification work under a trusted root and collision-resistant hashing. It gives no logarithmic bound on discovering every relevant event or on the human effort needed to understand it.

\subsubsection*{Proposition 5: Accountability Must Outlive the Process (\pdchapref{stp}{From Spawn to Person})}
Operating system processes are inherently mortal: they can be terminated by out-of-memory killers, abort on unhandled exceptions, or be restarted intentionally. If an agent's identity, reputation, debt ledger, and liabilities exist only in the volatile RAM of a running Python or Node process, a defecting agent can escape consequences simply by crashing and rebooting under a new process ID.

An Agent Operating Kernel decouples ephemeral execution bodies from durable agent personas. A task acquires a durable cryptographic principal (Ed25519 keypair), an append-only capability history, a Merkle-backed evidence record, active normative obligations, and ultimate liability for physical outcomes. Restarting a crashed process or leasing a fresh container does not wipe the case file clean or replenish consumed token budgets: a successor reconstructs recorded obligations from the durable ledger. Resuming the exact execution state additionally requires a valid execution checkpoint and a recovery protocol.

\begin{figure}[tp]
  \centering
  % Reader question: What survives when a process crashes?
% Claim: A durable principal can be bound to successive distinct process instances.
% Evidence: Principal A and two instance epochs.
% Distinguish: Principal, process identity and renewed authority.
% Grammar: Identity timeline, not automatic state resurrection.
% Acceptance: Identify why instance2 needs a fresh binding.
\pdfigurehue{pdcobalt}%
\begin{tikzpicture}[pd figure,x=1cm,y=1cm]

  \node[pd identity state,text width=8.6cm,align=center] (principal) at (5.2,0) {Durable principal A\\Key, obligations and retained receipts};
  \node[pd artifact,text width=3.7cm,align=center] (old) at (2.4,-2) {Instance 1\\PID 412; epoch 7};
  \node[pd focus state,text width=3.7cm,align=center] (new) at (8.0,-2) {Instance 2\\PID 919; epoch 8};
  \draw[pd arrow] (principal.south west) -- (old.north);
  \draw[pd arrow] (principal.south east) -- (new.north);
  \draw[pd breach arrow] (old) -- node[pd tag,text width=1.2cm,align=center] {crash} (new);
  \node[pd note,text width=9.2cm,align=left] at (5.2,-3.7) {The new instance proves a fresh binding and reacquires valid authority. Principal continuity does not revive expired leases or recover unrecorded state.};

\end{tikzpicture}
%
  \pdwidecaption{\textbf{Identity across process replacement.} A stable principal links successive execution instances. Epochs distinguish current authority from stale credentials, while committed obligations remain in durable storage.}[fig:spark-prop5]
\end{figure}

Formally, an agent persona is defined as a persistent state tuple:
\begin{equation}
\mathcal{P} \;=\; \langle \mathrm{pk}_{\text{agent}}, \mathcal{L}_{\text{history}}, \mathcal{O}_{\text{active}}, \mathcal{B}_{\text{consumed}}, \sigma_{\text{checkpoint}} \rangle
\end{equation}
where $\mathrm{pk}_{\text{agent}}$ is the public signing identity, $\mathcal{L}_{\text{history}}$ is the append-only hash chain of past commitments, $\mathcal{O}_{\text{active}}$ is the set of unfulfilled obligations, and $\mathcal{B}_{\text{consumed}}$ is the non-resetting spend tally. When a container process terminates ($\mathtt{SIGKILL}$), the kernel reclaims the ephemeral execution lease, but freezes $\mathcal{P}$. A successor must satisfy the current admission policy; the key alone does not override revocation, scope or epoch checks. The final component $\sigma_{\text{checkpoint}}$ is an optional execution checkpoint. Restoring obligations differs from reproducing exact inference state.

\subsubsection*{Proposition 6: Exchange Requires Separated Instruments (\pdchapref{he}{The Harbor Economy})}
When autonomous agents trade compute, task assignments, and data access, economic coordination requires robust market design. In classical economics, Akerlof~\pdcite{akerlof1970} demonstrated that asymmetric information and unseparated incentives degrade markets into a ``market for lemons.'' Conflating multiple financial functions into a single scalar reputation score or omnibus utility token creates acute moral hazard.

A trustworthy harbor economy strictly separates four distinct coordination instruments: (1)~\textbf{Payment}, which compensates delivered labor; (2)~\textbf{Collateral}, which bonds and prices bounded failure; (3)~\textbf{Evidence}, which authenticates specified completion artifacts whose semantic adequacy depends on the verifier; and (4)~\textbf{Settlement}, which atomically and irrevocably closes obligations. Combining these functions into a single score undermines trust and invites strategic manipulation; separating them makes their contracts and conservation checks independently inspectable. Audit quality still depends on the evidence and verifier.

\begin{figure}[tp]
  \centering
  % Reader question: which contract roles hold funds, support judgment and settle?
% Claim: distinct payment, collateral and evidence inputs feed settlement.
% Evidence: schematic contract roles, no empirical or unconditional incentive claim.
% Distinguish: buyer payment, worker liability, verifier evidence, final transfers.
% Grammar: aligned three-input accounting flow to one settlement stage.
% Acceptance: identify who funds each input and what permits the final transfers.
\begin{minipage}{11.1cm}
\centering\pdfigurehue{pdcobalt}%
\begin{tikzpicture}[pd figure,x=1cm,y=1cm]
  \def\cw{3.15}\def\gap{.45}\def\xzero{.05}
  \def\ytitle{7.90}\def\yparty{7.25}\def\yart{6.15}\def\yout{4.70}
  \pgfmathsetmacro{\xa}{\xzero+.5*\cw}
  \pgfmathsetmacro{\xb}{\xa+\cw+\gap}
  \pgfmathsetmacro{\xc}{\xb+\cw+\gap}
  \node[pd title,text width=3.1cm] at (\xa,\ytitle) {1. Payment};
  \node[pd title,text width=3.1cm] at (\xb,\ytitle) {2. Collateral};
  \node[pd title,text width=3.1cm] at (\xc,\ytitle) {3. Evidence};
  \node[pd label] at (\xa,\yparty) {buyer funds};
  \node[pd label] at (\xb,\yparty) {contractor funds};
  \node[pd label] at (\xc,\yparty) {contractor supplies};
  \node[pd value state,text width=2.65cm,minimum height=1.2cm] (pay)
    at (\xa,\yart) {tokens / fiat};
  \node[pd value state,text width=2.65cm,minimum height=1.2cm] (bond)
    at (\xb,\yart) {specified bond};
  \node[pd legible state,text width=2.65cm,minimum height=1.2cm] (ev)
    at (\xc,\yart) {hash chain and trace};
  \node[pd label,text width=2.9cm] (po) at (\xa,\yout) {held until accepted};
  \node[pd label,text width=2.9cm] (bo) at (\xb,\yout) {escrowed; liabilities stated};
  \node[pd label,text width=2.9cm] (eo) at (\xc,\yout) {checked against verifier terms};
  \draw[pd value arrow] (pay)--(po);
  \draw[pd value arrow] (bond)--(bo);
  \draw[pd legible arrow] (ev)--(eo);
  \node[pd climax state,text width=7.4cm,minimum height=1.45cm] (settle)
    at (\xb,1.70) {\textbf{4. Settlement}\\Release, refund or slash\\after the agreed verifier decision};
  \draw[pd value rule] (po.south)--(\xa,3.35)--(\xb,3.35);
  \draw[pd value rule] (bo.south)--(\xb,3.35);
  \draw[pd legible rule] (eo.south)--(\xc,3.35)--(\xb,3.35);
  \draw[pd focus arrow] (\xb,3.35)--(settle.north);
\end{tikzpicture}
\par\smallskip{\pdfiglabelitalic Separate contract roles expose accounting and verification
obligations. Incentive compatibility still depends on payoffs, detection and remedies.\par}
\end{minipage}
%
  \pdwidecaption{\textbf{Four contract roles.} Payment, collateral, evidence and settlement have distinct functions. Separating them exposes accounting and verification obligations; incentive compatibility additionally depends on payoffs, detection and enforceable remedies.}[fig:spark-prop6]
\end{figure}

Under this separated architecture, an agent offering a service must post collateral $C \ge \mathcal{V}(\text{risk})$ into escrow. When the task completes, the contractor produces evidence $\mathcal{E}$ (test logs, diff hashes). Only when $\mathcal{E}$ satisfies the verifier contract does settlement $\mathcal{S}$ execute, simultaneously disbursing payment $P$, returning collateral $C$, and updating the reputation ledger. Conflating payment and collateral into a single balance allows defaulting agents to consume resources without recourse; separating them exposes the quantities needed for an incentive analysis; it does not by itself make honest behavior optimal.

\subsubsection*{Proposition 7: Uncertainty Needs an Institution (\pdchapref{bonded}{The Bonded Commons})}
In complex software engineering, deterministic formal verification cannot decide every semantic property, and static type checkers cannot foresee every runtime interaction. Managing residual uncertainty requires institutional governance, grounded in Ostrom's governing of shared commons~\pdcite{ostrom1990governing} and Becker's economics of deterrence~\pdcite{becker1968}.

Some harms cannot be prevented upfront, and some outcomes cannot be judged purely by mechanical test suites. To govern this irreducible uncertainty, the kernel constructs institutional procedures: bonded stake deposits, randomized spot audits by independent witness verifiers, neutral contestation and appeal arbitration, and strict conservation rules. By turning uncertainty into an inspectable, contested process backed by financial slashing, the commons can deter modeled deviations when expected sanctions exceed gains, while providing a procedure for contesting false positives.

\begin{figure}[tp]
  \centering
  % Reader question: How can a contested outcome reach a terminal decision?
% Claim: A bond funds a specified procedure; deterrence needs an incentive model.
% Evidence: Task, audit, contest, appeal and settlement.
% Distinguish: Evidence evaluation and contractual finality.
% Grammar: A process with explicit branches, not an infallible justice device.
% Acceptance: Trace an appealed audit to a final decision.
\pdfigurehue{pdcobalt}%
\begin{tikzpicture}[pd figure,x=1cm,y=1cm]

  \node[pd value state,text width=5.6cm,align=center] (bond) at (5.2,0) {Agree the contract; post bond};
  \node[pd state,text width=5.6cm,align=center] (task) at (5.2,-1.4) {Perform work; retain evidence};
  \node[pd focus state,text width=5.6cm,align=center] (audit) at (5.2,-2.8) {Audit under a specified sampling rule};
  \node[pd ready state,text width=3.6cm,align=center] (accept) at (2.4,-4.4) {Uncontested outcome};
  \node[pd state,text width=3.6cm,align=center] (appeal) at (8.0,-4.4) {Contest and appeal\\Named decision-maker};
  \node[pd terminal,text width=5.6cm,align=center] (settle) at (5.2,-6.2) {Contractual finality and settlement};
  \draw[pd arrow] (bond) -- (task);
  \draw[pd arrow] (task) -- (audit);
  \draw[pd arrow] (audit) -- (accept);
  \draw[pd arrow] (audit) -- (appeal);
  \draw[pd arrow] (accept) -- (settle);
  \draw[pd arrow] (appeal) -- (settle);

\end{tikzpicture}
%
  \pdwidecaption{\textbf{Contestation and settlement.} Evidence enters a declared decision and appeal procedure. A final ruling authorizes settlement under the contract; its reliability depends on adjudication and enforcement assumptions.}[fig:spark-prop7]
\end{figure}

A one-shot deterrence calculation makes the assumptions explicit~\pdcite{becker1968}. Normalize honest utility to zero. Suppose a deviation gains $G\ge0$ when undetected, loses that gain and pays a slash $S\ge0$ when detected, and is detected with probability $p>0$. Then its expected incremental utility is
\begin{equation}
\Delta U=(1-p)G-pS.
\end{equation}
Thus $S\ge(1-p)G/p$ deters this deviation for a risk-neutral agent. For example, $G=10$, $p=0.1$, and $S=100$ give $\Delta U=-1$. Here $p$ includes both audit selection and successful detection. If the gain survives detection, the expression is $G-pS$ instead. A repeated-game equilibrium additionally needs a discount factor, monitoring model, continuation strategies and credible enforcement; this single-period inequality alone does not prove one.

\subsubsection*{Proposition 8: Federation Relays Evidence, Not Sovereignty (\pdchapref{fh}{The Federated Harbor})}
When autonomous fleets scale across distinct physical workstations, cloud tenancies, and organizational boundaries, a shared central database requires an explicit agreement about authority, availability and disclosure. This book studies hosts that retain independent admission authority.

The federation design here exchanges verifiable evidence while retaining local admission authority. Mutually distrustful hosts communicate across untrusted zero-trust relays, exchanging signed capability attestations, revocation gossips, Merkle inclusion proofs, and settlement receipts without sharing a database. Transport carries signed evidence under stated key-management assumptions; the receiving host's local kernel retains complete sovereignty over whether to admit, execute, or reject the requested operation.

\begin{figure}[tp]
  \centering
  % Reader question: Who decides whether a remote request may execute?
% Claim: The receiver verifies evidence under its own admission policy.
% Evidence: Two independently administered hosts and an untrusted relay.
% Distinguish: Transport, verification and local authority.
% Grammar: Trust-domain sequence, not shared-sovereignty imagery.
% Acceptance: Locate the receiving admission decision.
\pdfigurehue{pdcobalt}%
\begin{tikzpicture}[pd figure,x=1cm,y=1cm]

  \node[pd identity state,text width=6.0cm,align=center] (a) at (4.0,0) {Host A\\Signs request and supporting evidence};
  \node[pd protocol state,text width=6.0cm,align=center] (relay) at (4.0,-1.7) {Untrusted relay\\Carries signed bytes};
  \node[pd focus state,text width=6.0cm,align=center] (b) at (4.0,-3.4) {Host B admission check\\Trust, scope, freshness and local policy};
  \node[pd ready state,text width=3.5cm,align=center] (run) at (2.3,-5.2) {Admit locally\\Execute within local limits};
  \node[pd breach state,text width=3.5cm,align=center] (deny) at (7.7,-5.2) {Refuse locally\\Return a refusal receipt};
  \draw[pd arrow] (a) -- (relay);
  \draw[pd arrow] (relay) -- (b);
  \draw[pd ready arrow] (b) -- (run);
  \draw[pd breach arrow] (b) -- (deny);

\end{tikzpicture}
%
  \pdwidecaption{\textbf{Federated evidence.} A sender transmits signed evidence. The receiver checks identity, scope, freshness and local policy before admission. Authentication does not itself confer permission to act.}[fig:spark-prop8]
\end{figure}

In this federated topology, host kernels $H_A, H_B$ maintain private, independent SQLite WAL databases. When Agent~$\alpha$ on host $H_A$ requests execution on host $H_B$, it transmits a signed credential package $\Pi = \langle \mathrm{sig}_A(\alpha), \mathrm{cap}_{\alpha}, \mathcal{M}_{\text{proof}} \rangle$. Host $H_B$'s local kernel independently verifies $\Pi$ against its local trust anchor, checking revocation state without querying $H_A$'s internal database. Only if credential checks and local authorization both permit it may $H_B$ admit the task. Sovereignty remains local; evidence crosses the wire.

Read these as a dependency chain: each chapter uses earlier mechanisms and
states what they cannot guarantee.

\subsection{What Counts as Knowing: Epistemic Bounds and Evidence Channels}
\pdmarginfigure{lovelace}{Ada Lovelace observed that analytical engines have no pretensions to originate anything; they execute what we know how to order them to perform.}%
Every consequential statement in this volume has a status and a boundary. Derivations establish
relations under named assumptions. Finite checks exhaust a stated state space,
not every deployment. Simulations estimate behavior under a stated generator;
repository tests exercise named paths; running code proves availability, not
universal mediation. Counterexamples remain in the record because they mark
claims the research killed rather than conclusions it merely failed to prove.

Evidence is typed as well. When a chapter reports what a harbor observed, it
names the channel that produced the observation, and each channel has a blind
spot that no amount of volume removes:
\begin{itemize}[leftmargin=1.8em,itemsep=2pt]
  \item A \emph{deterministic verifier} (a test, a type checker, a model checker) shows that a stated property holds of the artifact it examined, and nothing about properties it did not state.
  \item A \emph{provenance witness} (a signature, a hash, a Merkle inclusion, an append-only log) shows attribution and tamper-evidence relative to a commitment, never that every event was logged or that a description is true.
  \item A \emph{protected meter} (a counter its subject cannot rewrite) shows a quantity, not its cause.
  \item \emph{Expert judgment} shows an assessment under a named rubric and carries the assessor's error rate.
  \item \emph{Buyer preference} shows what was chosen, not that the choice was right.
  \item A \emph{delayed outcome} shows what happened, and arrives too late to prevent it.
\end{itemize}
The reader should ask the same question of any dashboard: which channel, and what can it not show.

\Cref{app:volume-status} gives the implementation boundary;
\cref{app:result-atlas} records the results and their limits. The companion
\path{docs/harbor-research/mega-volume-epistemic-manifest.yaml} preserves each
result's source, status, conditions, boundary, chapter home, and verification
artifact. The chapters stand alone; the book's claim is the chain they make.

\section{Classical Multi-Agent Systems Foundations}\label{sec:prereq-classical-mas}

Before analyzing neural coding agents, we formalize the foundational frameworks developed during the classical era of multi-agent systems.

\subsection{Intentional Agency and Practical Rationality}

Classical decision theory models an agent as an expected utility maximizer balancing subjective beliefs $\mathcal{B}$ and desires $\mathcal{D}$. In his seminal 1987 work \emph{Intention, Plans, and Practical Reason}, philosopher Michael Bratman argued that intentions have a distinctive role in the practical reasoning of resource-bounded agents~\pdcite{bratman1987intention}.

A resource-bounded agent cannot continuously compute optimal policies from first principles at every millisecond. Practical rationality requires a third mental attitude: \textbf{Intention} ($\mathcal{I}$). Intentions are not merely desires with high utility; they represent \emph{conduct-controlling pro-attitudes to which the agent has actively committed}.

\pdmarginanalogy{prereq-flywheel-governor}{}{ch01-flywheel-governor}{Like a centrifugal flywheel governor throttling steam to stabilize rotational speed, practical rationality stabilizes an agent's intentions against transient sensor noise.}%
Bratman established four defining functional roles of intentions:
\begin{enumerate}[leftmargin=1.8em,itemsep=2pt]
  \item \textbf{Drive Means-End Reasoning:} An intention to achieve goal $\phi$ obligates the agent to formulate sub-plans and allocate computational cycles until $\phi$ is satisfied.
  \item \textbf{Filter of Admissibility:} Once an intention $\phi$ is adopted, candidate options inconsistent with $\phi$ (given current beliefs $\mathcal{B}$) are ruled inadmissible without evaluating their cost-benefit payoff.
  \item \textbf{Stability and Commitment:} Intentions resist reconsideration in the absence of relevant environmental interrupts or belief updates, preventing computational paralysis.
  \item \textbf{Coordination Scaffolding:} Intentions enable interpersonal coordination: other agents can form expectations, subject to reliability and enforcement assumptions, regarding an agent's future actions based on its declared commitments.
\end{enumerate}

\subsection{The Rao--Georgeff BDI Temporal Logic (\texorpdfstring{$BDI_{CTL^*}$}{BDI-CTL*})}

\paragraph{Example: A goal survives a failed attempt.}
Alice asks a worker to repair a failing test. The worker believes that revision~7
fails, wants a passing test, and commits to investigating the failure. Its first
patch does not compile. That observation refutes this attempted patch; it does
not establish that the repair is impossible. A different patch, a dependency
fix or a request for help may remain available. An intention tells us what
persists while those alternatives are considered.


Rao and Georgeff~\pdcite{rao1991modeling,rao1995bdi} formalize distinctions
between beliefs, selected goals and intentions in branching-time models.
A possible world represents a way circumstances could be; its branches represent
possible continuations. A belief modality quantifies over worlds the agent
regards as possible. It is an idealized model of reasoning, not a claim that
an interpreter stores every world or that a language model performs deduction.

For a small example let $W=\{w_7,w_8\}$, where the subscript is the repository
revision, and let $B_a\subseteq W$ be agent $a$'s nonempty set of possibilities.
Write $\operatorname{Bel}_a\phi$ when proposition $\phi$ holds in every world
in $B_a$. If $B_a=W$, the agent does not believe the revision is~7. It also does
not believe it is~8. If a trusted read narrows $B_a$ to $\{w_7\}$, it believes
revision~7 in this model. A stale read can make that belief false of the current
repository. This distinction between a modeled belief and an authoritative
observation returns in the confinement and supervision chapters.

\subsubsection{Compatibility and commitment}
An intention should remain subject to feasibility checks. Rao and Georgeff's
realism conditions relate belief, goal and intention worlds; their full
branching-time semantics are stronger and more specific than simply requiring
three lists to overlap. We use the practical distinction here without claiming
that the kernel implements their modal logic.

\subsection{Formal Commitment Strategies}
A commitment strategy states when reconsideration is permitted. The following
summaries follow the distinctions in Rao and Georgeff~\pdcite{rao1991modeling};
they are not transcriptions of the paper's axiom schemas.
\begin{table}[H]
\centering\small
\begin{tabularx}{\textwidth}{@{}lXX@{}}
\toprule
\textbf{Strategy} & \textbf{Persistence condition} & \textbf{Repair example}\\
\midrule
Blind & Persist until believed achieved. & Continue even after the required resource disappears.\\
Single-minded & Persist while achievement is still believed possible and not yet achieved. & Reconsider when that feasibility belief is lost.\\
Open-minded & Persistence also depends on retaining the goal. & Reconsider after Alice withdraws the repair request.\\
\bottomrule
\end{tabularx}
\caption{Commitment governs reconsideration, not the success of a plan.}
\label{tab:commitment-strategies}
\end{table}
Losing a belief that a repair is possible is not the same as believing it
impossible. In the two-world example, uncertainty can remove a belief without
establishing its negation. This matters when a tool error merely removes
confidence. The temporal ``until'' used for persistence in the cited model is
\emph{weak until}: persistence may continue forever if its release condition
never occurs. It therefore does not prove eventual completion.


\begin{pdclaim}{Design invariant}{Commitment Persistence}\label{inv:commitment-persistence}
An agent's recorded intention $\iota$ must not be abandoned merely because a single step incurs a transient tool failure. Intention reconsideration is an explicit state transition triggered by stated evidence, objective discharge, or authorized withdrawal or policy release. Context truncation is not itself such a transition.
\end{pdclaim}

\subsection{Operational Semantics of AgentSpeak(L) and Normative Extensions}

AgentSpeak(L)~\pdcite{bordini2007agentspeak} makes the distinction operational.
A belief base stores propositions, a plan library supplies triggered procedures,
and an intention stack records the remaining steps of an adopted plan. An event
selector chooses a pending event; an applicable-plan selector chooses a plan
whose trigger matches and whose context holds; an intention selector chooses
which stack advances. These are separate choices: having a recovery plan in
the library does not establish that it is applicable or scheduled.

For Alice's repair, ``test failed'' may trigger an investigation plan. A context
condition requiring a reachable test runner must still hold. If the runner is
down, selecting that plan and successfully running it are different questions.


Failure handling here follows the Jason variant of AgentSpeak~\pdcite{jasonfaq,jasonfailureimplementation}.
An empty applicable-plan set is a plan-selection failure, not evidence that a
host action was attempted and failed; reported action failure is a separate cause.
For an intention pursuing achievement goal $g$, Jason can seek a matching
$-!g$ handler in that goal's context or an enclosing goal's context. A handler
is usable only if its trigger matches and its context holds under the current
beliefs. Failure can propagate, and an intention without a usable handler may be
dropped. These cases follow separate interpreter paths.
Neither a failure event nor handler eligibility guarantees next-cycle execution
or eventual recovery, and earlier external effects are not automatically undone.

Figure~\ref{fig:bdi-execution-flow} shows how event, applicable-plan and intention
selection advance the BDI cycle, with a conditional return through goal-failure handling.

\begin{figure}[tp]
  \centering
  % Reader question: how does a BDI cycle select and recover an action?
% Claim: event, applicable-plan and intention selection advance the cycle.
% Evidence: schematic AgentSpeak(L) operational semantics, Bordini et al. 2007.
% Distinguish: belief filtering, selection functions, perceptions and failure events.
% Grammar: cycle with a belief artifact and a distinct failure return rail.
% Acceptance: trace a failure back to event selection without assuming success.
\begin{minipage}{11.1cm}
\centering\pdfigurehue{pdcobalt}%
\begin{tikzpicture}[pd figure,x=1cm,y=1cm]
  \def\dy{1.50}\def\sx{5.30}\def\gx{3.42}\def\lx{.65}
  \def\ax{9.05}\def\rx{10.70}
  \node[pd focus state,minimum width=3.2cm] (E) at (\sx,{4*\dy}) {event queue $\mathcal E$};
  \node[pd focus state,minimum width=3.2cm] (AP) at (\sx,{3*\dy}) {applicable plans};
  \node[pd focus state,minimum width=3.2cm] (IS) at (\sx,{2*\dy}) {intention stack};
  \node[pd focus state,minimum width=3.2cm] (XA) at (\sx,\dy) {execute $a$};
  \foreach \a/\b in {E/AP,AP/IS,IS/XA}{\draw[pd focus arrow] (\a)--(\b);}
  \foreach \i/\t/\n in {3.5/{select event}/1,2.5/{select plan}/2,1.5/{select intention}/3}{
    \node[pd label,anchor=east] at ({\gx-.32},{\i*\dy}) {\t};
    \node[pd badge] at (\gx,{\i*\dy}) {\n};}
  \node[pd artifact,text width=2.4cm] (BS) at (\ax,{3*\dy}) {belief base $\mathit{bs}$};
  \draw[pd thin arrow] (BS.west)--(AP.east);
  \draw[pd breach arrow,rounded corners=3pt]
    (XA.east)--(\rx,\dy)--(\rx,{4*\dy})--(E.east);
  \node[pd label,text=pderror,anchor=east] at ({\rx-.18},{3.5*\dy}) {$-!g$};
  \node[pd badge] at (\rx,3.0) {4};
  \node[pd label,text=pderror,text width=2.8cm,align=center]
    at (8.8,2.45) {goal failure; recovery may follow};
  \draw[pd focus arrow,rounded corners=3pt]
    (XA.west)--(\lx,\dy)--(\lx,{4*\dy})--(E.west);
  \node[pd focus label,rotate=90] at ({\lx-.40},{2.5*\dy}) {perceive};
\end{tikzpicture}
\par\smallskip
{\pdfiglabel\setlength{\tabcolsep}{3pt}%
\begin{tabularx}{\linewidth}{@{}lX@{}}
\toprule\textbf{Selection} & \textbf{What it does} \\\midrule
1 $\mathit{sel}_E(\mathcal E)$ & Choose an event: perception or an added goal. \\
2 $\mathit{sel}_O(\mathrm{App})$ & Choose an applicable plan; push it on the intention stack. \\
3 $\mathit{sel}_I(\mathit{is})$ & Choose the intention that advances by one action. \\
4 $-!g$ & A goal-context failure may raise a matching handler event; recovery needs a currently applicable plan and later selection. \\\bottomrule
\end{tabularx}}
\end{minipage}
%
  \pdwidecaption{\textbf{BDI selection flow.} Event, applicable-plan and intention selection determine the next action. Goal-failure handling needs a matching, applicable plan; recovery is not guaranteed. The diagram illustrates control flow, not successful execution.}[fig:bdi-execution-flow]
\end{figure}

\subsubsection{Normative Extensions: Deontic Constraints and Sanction Costs}

\begin{pdexample}{Two incompatible obligations}
Suppose one active norm requires retaining a diagnostic record and another
requires deleting that same record immediately. In this stipulated example,
violating retention costs 4 units and violating deletion costs 9. Assume the
only feasible choices satisfy exactly one norm. Retaining incurs 9; deleting
incurs 4. Minimizing the largest omitted sanction therefore chooses deletion.
This is a decision rule under supplied costs, not a proof of legal or ethical
priority. The chosen violation still needs an explanation and a record.
\end{pdexample}

In institutional multi-agent systems, agents are constrained by external norms (deontic logic). For this example, represent a norm by the six-tuple $N = \langle \mathrm{id}, \mathrm{Type}, \phi_{\text{act}}, \phi_{\text{tgt}}, d, \mathrm{Sanction} \rangle$, where $\mathrm{Type} \in \{\mathcal{O}, \mathcal{F}, \mathcal{P}\}$ (Obligation, Prohibition, Permission), $\phi_{\text{act}}$ is the activation context, $\phi_{\text{tgt}}$ is the target state, $d$ is the deadline, and $\mathrm{Sanction}$ is the penalty for violation.

An implementation may store general norms in an \textbf{Abstract Norm Base}
($ANB$) and instantiate them in a \textbf{Norm Instance Base} ($NIB$).
For a finite decision example, let $\mathcal N$ be the active norm set and let
$\emptyset\ne\Omega\subseteq 2^{\mathcal N}$ be a finite family of subsets that
are feasible under the stated action constraints and conflict-free under the
stated norm-conflict relation. Give each norm a stipulated finite nonnegative
sanction cost $c(n)$ for omission. One possible \textbf{worst-case sanction-cost}
rule is
\begin{equation}
S^*\in\arg\min_{S\in\Omega}
\max\bigl(\{c(n):n\in\mathcal N\setminus S\}\cup\{0\}\bigr).
\end{equation}
Here $\mathcal N$ is finite; including $0$ makes the maximum zero when no norm
is omitted. This minimizes the largest stipulated omitted sanction, not regret
relative to a counterfactual alternative. It supplies no audit record by itself.
If immutable evidence is required, a separate logging mechanism must commit the
chosen subset, any deliberate violations and their justification.

\subsubsection{Why BDI Matters for LLMs: Curing Intention Drift and Hallucinatory Abandonment}
\label{sec:why-bdi-matters}

The useful connection is a separation of jobs. A proposal generator can suggest
a new patch; a durable task record keeps track of what Alice still requested.
A compiler error can justify another attempt, a blocked state or an escalation.
It cannot by itself justify marking the obligation complete or proving that
no solution exists. Likewise, omitting the goal from a shortened prompt must
not silently delete it from the task record.

This is a design use of commitment theory, not evidence that a language model
has BDI mental states or that a particular kernel implements AgentSpeak.
Chapter 1 supplies the durable work-unit transitions; Chapter 5 asks which
obligations a replacement worker inherits. We return to the same failed repair
when considering timeouts: persistence also requires knowing whether an effect
already happened before attempting it again.


\subsection{Concurrency Foundations: Hewitt--Agha Actors vs. Hoare's Communicating Sequential Processes}

When multiple agents deliberate and execute concurrently on a developer's workstation, the concurrency substrate dictates whether the swarm operates harmoniously or degenerates into deadlocks and corrupted file state. Distributed systems research provides two foundational, mathematically rigorous concurrency paradigms:

\begin{itemize}[leftmargin=1.8em,itemsep=3pt]
  \item \textbf{Hewitt--Agha Actor Model (1973, 1986)~\pdcite{hewitt1973actor,agha1986actors}:}\enspace
  The Actor Model treats the \textbf{actor} as the universal primitive of concurrent computation. An actor is an autonomous, computational entity that encapsulates private state, a mailbox (message queue), and behavior. In response to a message, an actor can concurrently perform exactly three operations:
  \begin{enumerate}[leftmargin=1.5em,itemsep=1.5pt]
    \item \emph{Send} a finite number of messages to other actors with known addresses;
    \item \emph{Create} a finite number of new actors;
    \item \emph{Designate the replacement behavior} ($\mathtt{become}$) to be used for the next message it receives.
  \end{enumerate}
  Actors share \textbf{no mutable state}. Message passing is fundamentally \textbf{asynchronous}: sending is non-blocking, and messages are buffered in the recipient's mailbox.

  \item \textbf{Hoare's Communicating Sequential Processes (CSP) (1978)~\pdcite{hoare1978csp}:}\enspace
  CSP formalizes concurrent computation around \textbf{processes} communicating over synchronous, unbuffered \textbf{channels}. The core principles of CSP are:
  \begin{enumerate}[leftmargin=1.5em,itemsep=1.5pt]
    \item \emph{Synchronous Rendezvous:} Communication is an unbuffered handshake. The output event $c!v$ on channel $c$ and the input event $c?x$ constitute one joint communication event in the abstract semantics. Whichever process reaches the communication point first blocks until the other process arrives.
    \item \emph{Guarded Choice (Alternative):} A process can wait on multiple communication events simultaneously using the alternative construct ($\mathtt{alt}$), proceeding with whichever channel becomes ready first.
  \end{enumerate}
\end{itemize}

Figure~\ref{fig:actor-concurrency} details the dual concurrency architecture: Hewitt--Agha virtual actor isolation for stochastic deliberation, and Hoare's Communicating Sequential Processes (CSP) synchronous rendezvous for physical state mutations.

\begin{figure}[tp]
  \centering
  % Reader question: when may a sender continue relative to receiver readiness?
% Claim: buffered sends and synchronous rendezvous have different blocking rules.
% Evidence: two schematic schedules on a shared ordinal time axis.
% Distinguish: nonblocking send, receiver progress, joint event and optional timeout.
% Grammar: paired swimlanes with aligned time and a keyed semantics table.
% Acceptance: find the send that proceeds before its receiver is ready.
\begin{minipage}{11.1cm}
\centering\pdfigurehue{pdcobalt}%
\begin{tikzpicture}[pd figure,x=1cm,y=1cm]
  \def\xzero{2.05}\def\xe{9.85}\def\dy{1.2}\def\py{-4.2}
  \def\sendx{3.9}\def\recvx{7.8}
  \node[pd title,anchor=west] at (0,1.15) {Asynchronous actor send};
  \node[pd label,anchor=east,text width=1.75cm,align=right] at (1.75,0) {sender actor};
  \node[pd label,anchor=east,text width=1.75cm,align=right] at (1.75,-\dy) {receiver actor};
  \draw[pd focus arrow] (\xzero,0)--(\xe,0);
  \node[pd focus datum] at (\sendx,0) {};
  \node[pd label,anchor=east] at ({\sendx-.15},.48) {send};
  \node[pd badge] at ({\sendx+.22},.48) {1};
  \node[pd artifact,minimum width=2.6cm] (MB) at (5.85,-.6) {mailbox};
  \draw[pd thin arrow,rounded corners=2pt] (\sendx,-.08)--(\sendx,-.6)--(MB.west);
  \draw[pd thin arrow,rounded corners=2pt] (MB.east)--(\recvx,-.6)--(\recvx,{-\dy+.08});
  \draw[pd hairline] (\xzero,-\dy)--(\recvx,-\dy);
  \draw[pd focus arrow] (\recvx,-\dy)--(\xe,-\dy);
  \node[pd focus datum] at (\recvx,-\dy) {};
  \node[pd label,anchor=east] at ({\recvx-.15},{-\dy-.48}) {receive};
  \node[pd badge] at ({\recvx+.22},{-\dy-.48}) {2};
  \draw[pd thin arrow] (\xzero,{-\dy-.95})--(\xe,{-\dy-.95});
  \node[pd label,anchor=north east] at (\xe,{-\dy-1.05}) {time};
  \node[pd title,anchor=west] at (0,{\py+.7}) {Synchronous CSP rendezvous};
  \node[pd label,anchor=east] at (1.75,\py) {process $P_1$};
  \node[pd label,anchor=east] at (1.75,{\py-\dy}) {process $P_2$};
  \def\arrSx{4.0}\def\arrRx{5.35}\def\rvx{6.65}
  \foreach \y/\arr in {\py/\arrSx,{\py-\dy}/\arrRx}{
    \draw[pd protocol rule] (\xzero,\y)--(\arr,\y);
    \draw[pd hairline] (\arr,\y)--(\rvx,\y);
    \draw[pd protocol rule] (\rvx,\y)--(\xe,\y);}
  \draw[pd protocol rule] (\rvx,\py)--(\rvx,{\py-\dy});
  \node[pd protocol datum] at (\rvx,{\py-.5*\dy}) {};
  \node[pd badge] at ({\rvx-.55},{\py-.5*\dy}) {3};
  \node[pd protocol label,anchor=west,text width=2.85cm,align=left]
    at ({\rvx+.2},{\py-.5*\dy}) {joint communication};
  \node[pd breach state,minimum width=2.65cm] (TO) at (8.1,{\py-\dy-1.1}) {lease timeout};
  \draw[pd breach arrow,rounded corners=2pt]
    (\arrSx,\py)--(\arrSx,{\py-\dy-1.1})--(TO.west);
  \node[pd badge] at (4.55,{\py-\dy-1.45}) {4};
  \draw[pd thin arrow] (\xzero,{\py-\dy-1.95})--(\xe,{\py-\dy-1.95});
  \node[pd label,anchor=north east] at (\xe,{\py-\dy-2.05}) {time};
\end{tikzpicture}
\par\smallskip
{\pdfiglabel\setlength{\tabcolsep}{3pt}%
\begin{tabularx}{\linewidth}{@{}lX@{}}
\toprule\textbf{Event} & \textbf{Semantics} \\\midrule
1 Actor send & Sender proceeds without receiver readiness. \\
2 Actor receive & Receiver drains the mailbox and may adopt a new behaviour. \\
3 CSP rendezvous & One abstract joint event; the earlier arrival waits. \\
4 Timeout choice & An optional modeled guard offers another transition. \\\bottomrule
\end{tabularx}}
\par\smallskip{\pdfiglabelitalic Delays are illustrative. Ordering depends on the actor
implementation; rendezvous does not imply zero physical communication latency.\par}
\end{minipage}
%
  \pdwidecaption{\textbf{Two communication patterns.} Buffered task delivery separates sender and receiver timing. An admission exchange couples a request to its resource decision. Either actors or synchronous processes can implement safe resource ownership with the appropriate protocol.}[fig:actor-concurrency]
\end{figure}

\subsection{From Asynchronous Isolation to Synchronous Rendezvous}

The fundamental divergence between Hewitt's Actor model and Hoare's Communicating Sequential Processes reflects an architectural trade-off between \emph{isolation} and \emph{synchronization}:

\begin{enumerate}[leftmargin=1.8em,itemsep=3pt]
  \item \textbf{Asynchronous Mailboxes vs.\ Unbuffered Barriers:}\enspace In the Actor model, communication is non-blocking and decoupled in time. An agent posts a message to a peer's mailbox and immediately resumes local execution. This temporal decoupling helps keep a supervisory loop responsive while workers perform long-running inference. A synchronous design can achieve similar responsiveness by separating those waits into worker processes. However, unbounded asynchronous mailboxes introduce severe hazards: message ordering is non-deterministic, queue backpressure is difficult to meter, and an agent can be flooded with stale context deltas.
  
  \item \textbf{Synchronous Rendezvous for Scarce Resources:}\enspace An unbuffered send waits for a corresponding receive. This can express an admission handshake, but does not itself establish mutual exclusion, durability or deadlock freedom. Those properties depend on the protocol surrounding the handshake.

  \item \textbf{The Synthesis in an Agent Operating Kernel:}\enspace Buffered task delivery separates deliberation from resource admission. A resource manager can serialize conflicting operations, validate versions and issue fencing tokens. It may itself be implemented as an actor or a synchronous process; correctness is a property of its transition rules and actuator, not of the concurrency model's name.
\end{enumerate}

\subsubsection{Why Concurrency Lineages Matter: Protocol Obligations}
\label{sec:why-concurrency-matters}

An actor can safely own a shared resource and process requests in order. A CSP network can avoid deadlock through its protocol structure and suitable timeout or cancellation behavior. Conversely, either model can contain an erroneous resource protocol. In the design studied here, workers deliberate asynchronously and request authority at an explicit admission boundary. Leases need fencing so that an expired holder cannot later act; a timeout by itself does not prevent a stale write.

\subsection{Communicative Action and Speech Acts (FIPA-00037)}

A message has a role in a conversation as well as a byte representation. \textbf{Speech-act theory} distinguishes what is said, the act performed in saying it, and the consequences of saying it. Austin distinguishes locutionary acts (meaningful utterances), illocutionary acts (such as requests or promises), and perlocutionary acts (such as persuading a listener)~\pdcite{austin1962how}. Searle develops a systematic account of illocutionary acts~\pdcite{searle1969speech}. The Foundation for Intelligent Physical Agents standardized the \textbf{FIPA ACL Communicative Act Library} (FIPA-00037)~\pdcite{fipa2002acl}.

Consider an operator saying, ``Please test revision 7.'' The words identify a
particular requested action: that is their meaning. In this setting the
operator is making a request, rather than reporting a completed test: that is
its force. The request may persuade the worker to stop another task and run
the test: that is a possible consequence. A worker can understand the request
and still refuse it. Recording the utterance therefore does not establish
that the requested work occurred, or that an obligation was accepted.

FIPA distinguishes the sender's feasibility conditions from the intended
rational effect of a communicative act. An intended effect is not a guarantee
that a receiver believes, agrees or acts. For our repair workflow, the following
illustrative message kinds keep different protocol events apart. This table
is not a reproduction of FIPA's formal semantics.

\begin{table}[H]
\centering\small
\begin{tabularx}{\textwidth}{@{}l>{\raggedright\arraybackslash}X>{\raggedright\arraybackslash}X@{}}
\toprule
\textbf{Kind} & \textbf{Example content} & \textbf{What remains to check}\\
\midrule
\texttt{request} & Run the test on revision 7. & Authorization and acceptance.\\
\texttt{inform} & The test passed. & Which run, revision and evidence?\\
\texttt{propose} & I can do it for 20 units. & Terms and explicit acceptance.\\
\texttt{agree} & I accept job J under those terms. & Whether the protocol records a binding acceptance.\\
\texttt{refuse} & I cannot accept job J. & Whether another worker is needed.\\
\texttt{failure} & Attempt A did not complete. & Which effects, if any, already occurred?\\
\bottomrule
\end{tabularx}
\caption{Illustrative workflow messages inspired by communicative-act distinctions.}
\label{tab:fipa-performatives}
\end{table}


\begin{figure}[tp]
  \centering
  % Reader question: which responses may follow a communicative act?
% Claim: request, inform and propose express intent, not physical enforcement.
% Evidence: upstream protocol lifecycles with named performatives and outcomes.
% Distinguish: an assertion, its assessment and a possible belief revision.
% Grammar: preserve upstream small-multiple state machines.
% Acceptance: an inform message does not itself verify its asserted fact.
\pdfigurehue{pdcobalt}%
\begin{tikzpicture}[pd figure,x=1cm,y=1cm]
  % Native 10.9cm grid; the three protocols share columns and row pitch.
  \def\colA{1.55}\def\colB{5.10}\def\colC{9.05}
  \def\rowdy{3.20}\def\branch{0.65}
  \foreach \i/\act/\kind in {0/request/directive,1/inform/assertive,2/propose/commissive}{
    \node[pd title,anchor=west] at (0.05,{-\i*\rowdy+1.20})
      {\number\numexpr\i+1\relax. \texttt{\act} --- \kind};}
  \node[pd state,minimum width=2.2cm] (R1) at (\colA,0) {\texttt{request}};
  \node[pd state,minimum width=2.0cm] (R2) at (\colB,0) {\texttt{agree}};
  \node[pd ready state,minimum width=3.0cm] (R3) at (\colC,\branch) {\texttt{inform-done}};
  \node[pd breach state,minimum width=2.4cm] (R4) at (\colC,-\branch) {\texttt{failure}};
  \node[pd breach state,minimum width=2.2cm] (RX) at (\colA,-1.15) {\texttt{refuse}};
  \draw[pd initial] (0.05,0) -- (R1.west);
  \draw[pd arrow] (R1) -- (R2);
  \draw[pd ready arrow] (R2) -- (R3);
  \draw[pd breach arrow] (R2) -- (R4);
  \draw[pd breach arrow] (R1) -- (RX);
  \node[pd state,minimum width=2.2cm] (I1) at (\colA,-\rowdy) {\texttt{inform}($\phi$)};
  \node[pd state,text width=2.45cm] (I2) at (\colB,-\rowdy) {assess assertion};
  \node[pd ready state,text width=2.7cm] (I3) at (\colC,-\rowdy) {possible revision};
  \draw[pd initial] (0.05,-\rowdy) -- (I1.west);
  \draw[pd arrow] (I1) -- (I2);
  \draw[pd ready arrow] (I2) -- (I3);
  \node[pd state,minimum width=2.2cm] (P1) at (\colA,{-2*\rowdy}) {\texttt{propose}};
  \node[pd state,minimum width=2.2cm] (P2) at (\colB,{-2*\rowdy}) {utility eval};
  \node[pd ready state,minimum width=3.2cm] (P3) at (\colC,{-2*\rowdy+\branch}) {\texttt{accept-proposal}};
  \node[pd breach state,minimum width=3.2cm] (P4) at (\colC,{-2*\rowdy-\branch}) {\texttt{reject-proposal}};
  \draw[pd initial] (0.05,{-2*\rowdy}) -- (P1.west);
  \draw[pd arrow] (P1) -- (P2);
  \draw[pd ready arrow] (P2) -- (P3);
  \draw[pd breach arrow] (P2) -- (P4);
  \node[pd title,anchor=west] at (.05,-8.00) {Outcome key};
  \foreach \i/\s/\w in {0/{pd state}/{performative or internal assessment},
    1/{pd ready state}/{protocol outcome; not proof of content},
    2/{pd breach state}/{refusal or failure}}{
    \node[\s,minimum width=.55cm,minimum height=.36cm,inner sep=0pt] at (.35,{-8.65-\i*.55}) {};
    \node[pd label,anchor=west] at (.85,{-8.65-\i*.55}) {\w};}
\end{tikzpicture}
%
  \pdwidecaption{\textbf{Message intent and protocol state.} A request asks for action, an inform message asserts a proposition, and a proposal offers terms. Authentication identifies the sender; it does not make an assertion true or a proposal binding.}[fig:fipa-speech-acts]
\end{figure}

\subsubsection{Why Speech Acts Matter: From Conversational Fluff to Cryptographic Envelopes}
\label{sec:why-fipa-matters}

``I will look into it'' might acknowledge receipt, express willingness or accept
an obligation. A protocol can remove this ambiguity by requiring a message kind,
job identifier, sender, recipient and explicit terms. A signature authenticates
those bytes under a key; it does not make ``the test passed'' true. An
\texttt{agree} message becomes an acceptance only if the governing protocol
admits and records it.

Schema validation and authentication therefore help with distinct checks. They
do not remove malicious content from a message body or prove compliance with a
promise. Chapter 2 develops authority, and Chapter 6 separates a proposal,
acceptance, delivery evidence and settlement. These are design requirements,
not a claim that every existing messaging endpoint implements this envelope.


\subsection{Task Allocation: Smith's Contract Net Protocol (1980)}

Reid G. Smith's (1980) \textbf{Contract Net Protocol (CNP)}~\pdcite{smith1980cnp} formalized bilateral market-inspired negotiation between a Manager and distributed Contractor agents. A possible coding-workflow adaptation has four stages:

\begin{enumerate}[leftmargin=1.8em,itemsep=2pt]
  \item \textbf{Call for Proposals (CFP Broadcast):}\enspace The Manager broadcasts an announcement $\langle\mathtt{task\_id}, \allowbreak \mathtt{eligibility}, \allowbreak \mathtt{spec}, \allowbreak t_{\text{expire}}\rangle$ specifying required toolsets, maximum spend budget, and deadline.
  \item \textbf{Bidding and Capability Self-Assessment:}\enspace Eligible contractor agents inspect their current local capacity, sandboxed tools, and model context before submitting formal bids $\langle\mathtt{bid\_id}, \allowbreak \mathtt{cost}, \allowbreak \mathtt{eta}, \allowbreak \mathtt{reputation}\rangle$, or returning explicit refusals if overloaded.
  \item \textbf{Bid Evaluation and Award:}\enspace The manager selects among eligible bids using a declared rule. For example, among bids meeting the same deadline and acceptance specification, select the lowest price. Resource reservation requires a separate admission step.
  \item \textbf{Execution, Interim Streaming, and Settlement:}\enspace The contractor checks out an isolated worktree branch, streams structured progress deltas, commits the diff, runs verification tests, and submits cryptographic execution receipts to unlock escrow payment.
\end{enumerate}

\begin{figure}[tp]
  \centering
  % Reader question: who sends what before allocation, execution and settlement?
% Claim: Contract Net allocation precedes local admission and verifier-gated settlement.
% Evidence: two contractors, one manager, six numbered message types.
% Distinguish: proposal, refusal, contractual award, host admission and settlement.
% Grammar: preserve the upstream message ladder and ruled payload legend.
% Acceptance: locate the award before the admitted execution receipt.
\noindent\begin{minipage}{11.1cm}\centering
\pdfigurehue{pdcobalt}%
\begin{tikzpicture}[pd figure,x=1cm,y=1cm]
  % a comma-joined style list cannot live inside a \foreach macro field (the
  % node-options parser splits on RAW commas before macro expansion, so a
  % macro holding "roleA,roleB" is read back as one literal unknown key); a
  % named compound style is the fix, composed only from existing pd roles.
  \tikzset{
    pd breach badge/.style={pd badge,pd breach fill},
    pd legible badge/.style={pd badge,pd legible fill},
    pd value badge/.style={pd badge,pd value fill},
    pd value focus arrow/.style={pd value rule,->},
  }
  % ---- one grid, declared once ---------------------------------------------
  \def\dy{0.80}    % message pitch
  \def\gap{4.10}   % lifeline spacing: contractor 1 | manager | contractor 2
  \def\yo{-1.05}   % centre line of message 1
  % ---- heads and lifelines -------------------------------------------------
  \foreach \i/\name in {0/Contractor 1, 1/Manager, 2/Contractor 2}{
    \node[pd state,minimum width=2.7cm] (h\i) at ({1.35+\i*\gap},0) {\name};
    \draw[pd guide] ({1.35+\i*\gap},-0.52) -- ({1.35+\i*\gap},{\yo-5*\dy-0.50});}
  % ---- messages: from, to, number, arrow, badge -----------------------------
  % Every arrow slopes forward in time (tail above head).  Hue is meaning:
  % ink = ordinary traffic, focus = the award, red dashed = the refusal,
  % teal = evidence, gold = the only payment.  What each numbered message IS
  % lives once, in the ruled legend below -- not repeated as a word on the
  % arrow too (S7: a relation word beside its edge, never stacked on a badge
  % that already carries a legend entry).
  \foreach \a/\b/\k/\s/\bs in {%
    1/0/1/{pd arrow}/{pd badge},
    1/2/1/{pd arrow}/{pd badge},
    0/1/2/{pd arrow}/{pd badge},
    2/1/3/{pd breach arrow}/{pd breach badge},
    1/0/4/{pd focus arrow}/{pd focus badge},
    0/1/5/{pd legible arrow}/{pd legible badge},
    1/0/6/{pd value focus arrow}/{pd value badge}}{
      \pgfmathsetmacro{\yk}{\yo-(\k-1)*\dy}
      \pgfmathsetmacro{\xm}{1.35+(\a+\b)*\gap/2}
      \node[\bs] (m\a\b\k) at ({\xm},{\yk}) {\k};
      \draw[\s,-,pd to badge]  ({1.35+\a*\gap},{\yk+0.12}) -- (m\a\b\k);
      \draw[\s,pd from badge]  (m\a\b\k) -- ({1.35+\b*\gap},{\yk-0.12});}
\end{tikzpicture}
\par\vspace{4pt}
\pdfiglabel
\begin{tabularx}{\linewidth}{@{}c >{\raggedright\arraybackslash}p{2.1cm} >{\raggedright\arraybackslash}X@{}}
\toprule
\# & message & payload \\
\midrule
1 & call for proposals & task \texttt{build-kernel-ffi}; spend cap \$0.00; tools \texttt{cargo}, \texttt{git}, \texttt{seatbelt} \\
2 & bid & eta 12\,min; cost \$0.00; reputation 0.98 \\
3 & refusal & lock contention --- the contractor is busy \\
4 & award & explicit contract terms and exclusive worktree lease \\
5 & receipts & admitted execution result, diff and verifier test run \\
6 & settlement & release, refund or slash under agreed verifier terms \\
\bottomrule
\end{tabularx}
\par\smallskip
{\pdfiglabelitalic\raggedright Local authority is checked before effects; an award does not bypass admission.\par}
\end{minipage}
%
  \pdwidecaption{\textbf{Contract-net task allocation.} A manager announces a task, evaluates bids and awards work. A contractor returns evidence for contract-specific acceptance and settlement. Allocation, execution authority and payment verification are separate obligations.}[fig:cnp-protocol]
\end{figure}

\vspace{10pt}
\noindent\begin{minipage}{\linewidth}
  \centering
  \pdswissfigplate{fig-ch0-cnp-bidding}{0.88\linewidth}%
  \par\vspace{8pt}%
  {\color{pdcobalt}\rule{\linewidth}{1.2pt}}\par\vspace{5pt}%
  \raggedright
  {\color{pdswissred}\rule{4.5pt}{4.5pt}}\enspace%
  {\pdgrotesk\bfseries\scshape\fontsize{9}{11}\selectfont\color{hhink}%
    Architectural Plate 0.A\enspace\textperiodcentered\enspace Task Allocation Dynamics (Contract Net Protocol)}\par\vspace{4pt}%
  {\pdgrotesk\itshape\fontsize{9.5}{13}\selectfont\color{hhink}%
    ``Workers assess local capacity; a manager selects a bid. Resource authority and acceptance evidence remain separate checks.''}\par\vspace{6pt}%
  {\color{hhgray!30}\rule{\linewidth}{0.4pt}}
\end{minipage}
\vspace{10pt}

\subsubsection{Why CNP Matters: Selection Is Not Acceptance}
\label{sec:why-cnp-matters}

In any multi-agent engineering swarm, allocating tasks across workers presents a fundamental distributed systems dilemma:
\begin{enumerate}[leftmargin=1.8em,itemsep=2pt]
  \item \textbf{The Bottleneck of Central Planning:} If a centralized orchestrator agent must manually inspect each worker's context window, parse their local git branches, and micromanage work assignment, the orchestrator becomes a severe throughput bottleneck and single point of cognitive failure.
  \item \textbf{The Chaos of Greedy Contention:} Conversely, if workers greedily claim tasks from an uncoordinated queue, multiple agents race to modify the same files, squandering token budgets on overlapping AST diffs and inducing git merge conflicts.
\end{enumerate}

\begin{pdexample}{A cheaper bid can still fail acceptance}
Two eligible workers offer to test revision~7 before the same deadline, for
20 and 24 units. The lowest-price rule selects the first. It returns a passing
receipt for revision~6. The selection rule has worked; the delivery has not
met the agreed \mbox{specification}. Award, permission to execute and accepted delivery
are three different events. Chapter 6 develops their ledger consequences.
\end{pdexample}

Reid Smith's Contract Net Protocol (CNP)~\pdcite{smith1980cnp} supplies a
negotiation pattern for allocating tasks. Applying it to a coding workflow
requires separate contracts for resource ownership and settlement:
\begin{itemize}[leftmargin=1.8em,itemsep=2pt]
  \item \textbf{Proposed task-allocation ceremony.} Broadcast a task and budget;
  receive capability statements and bids; select a bid under a declared rule;
  reserve the required resources; and settle against specified delivery evidence.
  Bidding alone does not establish that the winner can safely execute the work.
  The checked-in dispatch runner plans queued work for a selected backend,
  budget, and isolated worktree. The Harbor Work Register records cooperative
  claims and refuses a second live holder. Those components do not by themselves
  implement competitive bidding or escrow settlement.
  \item \textbf{Confinement and evidence obligations.} Execution confinement and
  spend controls depend on the selected execution path. A settlement contract
  may require authenticated test receipts and commit digests, but these bind
  an identified run and artifact to a signer; they do not prove every semantic
  property of the work. Resource reservation, verifier trust, payment authority,
  and the conditions for refusing settlement must each be specified.
\end{itemize}

\vspace{10pt}
\noindent\begin{minipage}{\linewidth}
  \centering
  \pdswissfigplate{fig-ch0-cnp-settlement}{0.88\linewidth}%
  \par\vspace{8pt}%
  {\color{pdcobalt}\rule{\linewidth}{1.2pt}}\par\vspace{5pt}%
  \raggedright
  {\color{pdswissred}\rule{4.5pt}{4.5pt}}\enspace%
  {\pdgrotesk\bfseries\scshape\fontsize{9}{11}\selectfont\color{hhink}%
    Architectural Plate 0.B\enspace\textperiodcentered\enspace Bilateral Escrow Settlement}\par\vspace{4pt}%
  {\pdgrotesk\itshape\fontsize{9.5}{13}\selectfont\color{hhink}%
    ``A settlement contract can require signed test receipts and commit digests; those receipts attest to a run, not to every property of the delivered software.''}\par\vspace{6pt}%
  {\color{hhgray!30}\rule{\linewidth}{0.4pt}}
\end{minipage}
\vspace{10pt}

\subsection{Dynamic Epistemic Logic and Big Brother Logic}

\begin{pdexample}{A hidden bit and an observed effect}
A private file contains a bit, either 0 or 1. A worker sees the same message
``ready'' in both cases. Relative to this observation model, it cannot tell
which bit is stored. Now suppose the worker also sees an exit code equal to
the bit. Hiding the file no longer hides the fact. Confinement must account for
both observations. Chapter 3 develops this problem for sequences of releases.
\end{pdexample}

Let $w_0,w_1$ denote these two states. Write $w R_i w'$ when agent $i$ cannot
distinguish the states using its allowed observations. The assertion $K_i\phi$
means that proposition $\phi$ holds in every state accessible from the current
one by $R_i$. Before the exit-code observation, both states are accessible;
after it, only the matching state remains. Unlike the belief example above,
a knowledge model includes the actual state among its possibilities.

Charrier et al.~\pdcite{charrier2014bigbrother} study an observation model
called Big Brother Logic, using cameras and fields of view. Its geometrical
restrictions make asymmetric visibility concrete. The engineering lesson is
the need to specify observations before asserting knowledge or secrecy.
Seeing another agent is not the same as seeing all that agent's evidence.


\begin{figure}[tp]
  \centering
  % Reader question: when can observing a worker support higher-order knowledge?
% Claim: two directed observations do not imply transitive direct visibility.
% Evidence: a2 lies in V(a1); phi lies in V(a2), outside V(a1).
% Distinguish: visible objects, known relations and higher-order knowledge.
% Grammar: upstream cones and arrows with a conditional legend.
% Acceptance: identify the extra observation-semantics assumption.
\pdfigurehue{pdcobalt}%
\begin{tikzpicture}[pd figure,x=1cm,y=1cm]
  % ---- the grid -------------------------------------------------------------
  \def\Wl{0.0}  \def\Wr{10.7}  \def\Hb{0.40}  \def\Ht{5.40}
  \coordinate (a1) at (0.75,2.00);
  \coordinate (a2) at (4.45,2.90);
  \coordinate (phi) at (8.05,3.80);
  \def\angA{13.2}  \def\spreadA{16}  \def\radA{6.60}   % Vis(i): wide, long
  \def\angB{12.7}  \def\spreadB{17}  \def\radB{4.60}  % Vis(j): narrower
  \def\trim{0.52}  % clears the 8.5mm agent node, so no edition's pattern
                    % fill ever needs to occlude a stroke under its own text
  % ---- the field: unobserved by default, named once -----------------------
  \draw[pd neutral fill] (\Wl,\Hb) rectangle (\Wr,\Ht);
  \node[pd label,anchor=south west,text=pdinkmuted] at (\Wl+0.18,\Hb+0.14)
    {epistemic shadow};
  % ---- V(a1): wide, reaches far enough to hold a2 --------------------------
  % the wedge's near-apex edge is trimmed clear of the agent node (see
  % \trim above) rather than meeting at the node's exact centre, so the
  % stroke never needs the node's own fill to occlude it.
  \fill[pd focus fill]
    ($(a1)+({\angA+\spreadA}:\trim)$) -- ($(a1)+({\angA+\spreadA}:\radA)$)
      arc[start angle={\angA+\spreadA},end angle={\angA-\spreadA},radius=\radA]
      -- ($(a1)+({\angA-\spreadA}:\trim)$) -- cycle;
  \node[pd focus label,anchor=south west] at ($(a1)+(1.2,1.55)$) {$\mathrm{Vis}(i)$};
  % ---- V(a2): narrower, reaches far enough to hold phi ---------------------
  \fill[pd legible fill]
    ($(a2)+({\angB+\spreadB}:\trim)$) -- ($(a2)+({\angB+\spreadB}:\radB)$)
      arc[start angle={\angB+\spreadB},end angle={\angB-\spreadB},radius=\radB]
      -- ($(a2)+({\angB-\spreadB}:\trim)$) -- cycle;
  \node[pd legible label,anchor=south west] at ($(a2)+(0.55,1.35)$) {$\mathrm{Vis}(j)$};
  % ---- the two agents and the objective, named nodes so every incident
  % line clips to the node's own boundary -- never through its text, in any
  % edition's fill treatment (hatched editions do not occlude an overlaid
  % line the way a solid tint does).
  \node[pd automaton focus,minimum size=8.5mm] (node-a1) at (a1) {$i$};
  \node[pd automaton state,minimum size=8.5mm] (node-a2) at (a2) {$j$};
  \node[pd datum] (node-phi) at (phi) {};
  \node[pd label,anchor=west] at ($(phi)+(0.28,0)$) {$\phi$};
  % ---- the two observation relations, each an arrow with a badge ----------
  \draw[pd focus arrow] (node-a1) -- (node-a2);
  \draw[pd legible arrow] (node-a2) -- (node-phi);
  \node[pd badge] at ($(a1)!0.58!(a2)$) {1};
  \node[pd badge] at ($(a2)!0.58!(phi)$) {2};
  % ---- the ruled legend: the two observations compose into the verdict ----
  % badges are set as bold numerals, not redrawn circles, so the table never
  % depends on font inheritance into a nested tikzpicture.
  \node[anchor=north west,text width=10.7cm,align=left] at (\Wl,\Hb-0.35) {%
    \renewcommand{\arraystretch}{1.28}%
    \pdfiglabel
    \begin{tabular}{@{}cp{3.55cm}p{4.65cm}@{}}
      \toprule
      & \textbf{condition} & \textbf{reading} \\
      \midrule
      \textbf{1} & $j \in \mathrm{Vis}(i)$ & $i$ observes $j$ \\
      \textbf{2} & $\phi \in \mathrm{Vis}(j)$ & $j$ observes $\phi$ \\
      $1 \wedge 2$ & known $\mathrm{Vis}(j)$ and observation semantics & may yield $\mathsf{K}_i\mathsf{K}_j\phi$; no direct visibility or covert-channel guarantee \\
      \bottomrule
    \end{tabular}};
\end{tikzpicture}
%
  \pdwidecaption{\textbf{Asymmetric visibility.} The observation graph records who sees which modeled object. Visibility is not automatically transitive: seeing an agent does not establish knowledge of everything it knows. Host confidentiality needs an additional channel model.}[fig:big-brother-logic]
\end{figure}

\vspace{10pt}
\noindent\begin{minipage}{\linewidth}
  \centering
  \pdswissfigplate{fig-ch0-big-brother-logic}{0.88\linewidth}%
  \par\vspace{8pt}%
  {\color{pdcobalt}\rule{\linewidth}{1.2pt}}\par\vspace{5pt}%
  \raggedright
  {\color{pdswissred}\rule{4.5pt}{4.5pt}}\enspace%
  {\pdgrotesk\bfseries\scshape\fontsize{9}{11}\selectfont\color{hhink}%
    Architectural Plate 0.C\enspace\textperiodcentered\enspace Dynamic Epistemic Cones \& Confinement}\par\vspace{4pt}%
  {\pdgrotesk\itshape\fontsize{9.5}{13}\selectfont\color{hhink}%
    ``Observation models specify who can distinguish which states; confinement must justify the correspondence to actual host channels.''}\par\vspace{6pt}%
  {\color{hhgray!30}\rule{\linewidth}{0.4pt}}
\end{minipage}
\vspace{10pt}

\subsubsection{Why Big Brother Logic Matters: Asymmetric Observation and Zero-Trust Confinement}
\label{sec:why-bbl-matters}

In multi-agent systems research, information sharing is often naively treated as an unalloyed good: the more agents communicate, the better they coordinate. In production systems engineering, this assumption is catastrophic. Unrestricted visibility leaks proprietary credentials, inflates context windows with irrelevant chatter, and enables prompt-injection side-channels.

Observation models help separate three engineering obligations:
\begin{itemize}[leftmargin=1.8em,itemsep=2pt]
  \item \textbf{Recorded effects:} A supervisor needs authentic receipts for the events it claims to observe. Recording a tool-call request is not a complete observation of every effect of its subprocess.
  \item \textbf{Indistinguishability:} To model secrecy of a proposition $\phi$ at a world, require accessible alternatives with $\phi$ true and false, yielding $\neg K_i\phi\land\neg K_i\neg\phi$. Hidden files alone do not prove this if another observation reveals the same fact.
  \item \textbf{Confinement:} A sandbox controls specified host channels. Proving noninterference requires relating those channels and observations to the epistemic model. BBL's camera geometry is an explanatory model, not a proof of operating-system confidentiality.
\end{itemize}

\subsection{Game-Theoretic Foundations and Mechanism Design}

Return to the deterrence example: gain 10, detection probability 0.1 and
sanction 100 give deviation payoff $0.9(10)-0.1(100)=-1$. If detection falls
to 0.01, that payoff becomes $0.99(10)-0.01(100)=8.9$. The same bond no longer
deters the deviation. Incentives depend on the information and enforcement
model, not just the amount displayed in an account. Chapters 5--7 repeatedly
use this distinction.

When agents have independent objective functions, coordination cannot assume altruism. Shoham and Leyton-Brown~\pdcite{shoham2009multiagent} systematized the algorithmic and game-theoretic foundations of multi-agent systems:

\begin{enumerate}[leftmargin=1.8em,itemsep=3pt]
  \item \textbf{Normal-Form Games and Nash Equilibrium:}\enspace Let $N=\{1,\ldots,n\}$ index players, $A_i$ be player $i$'s action set, $S_i$ its mixed strategies (probability distributions over actions), and $u_i$ its expected payoff. Write $s_{-i}$ for the other players' strategies. A strategy profile $s^* = (s_1^*, \dots, s_n^*)$ is a Nash equilibrium if no agent can unilaterally deviate and increase its expected payoff:
  \begin{equation}
  \forall i \in N, \quad \forall s_i \in S_i, \quad u_i(s_i^*, s_{-i}^*) \;\ge\; u_i(s_i, s_{-i}^*)
  \end{equation}

  \item \textbf{Mechanism Design and the Revelation Principle:}\enspace A mechanism specifies allocations and transfers. The revelation principle says that an outcome implementable in dominant strategies, or in Bayes--Nash equilibrium, has a direct mechanism with truthful reporting in the corresponding equilibrium sense. It does not say that every desired outcome is implementable. Here $\theta_i$ is a private type, $\hat\theta_i$ a report, $k$ the allocation rule, $v_i$ the value of an allocation to type $\theta_i$, and $p_i$ the payment charged to $i$. Quasi-linear utility is value minus payment. Dominant-strategy truthfulness requires, for every $i$, $\theta_i$, $\hat\theta_i$ and $\hat\theta_{-i}$,
  \begin{align}
  v_i(k(\theta_i,\hat\theta_{-i});\theta_i)-p_i(\theta_i,\hat\theta_{-i})
  &\ge v_i(k(\hat\theta_i,\hat\theta_{-i});\theta_i)\nonumber\\
  &\quad-p_i(\hat\theta_i,\hat\theta_{-i}).
  \end{align}
  Bayesian incentive compatibility instead compares the corresponding expected utilities over other types under the stated prior.

  \item \textbf{Vickrey--Clarke--Groves (VCG) Mechanisms:}\enspace For intuition, auction one indivisible slot. A bidder values it at 10,
and the highest rival bid is 5. Winning and paying 5 gives utility 5; losing
gives 0. Any bid above 5 wins at that same price, so truthfulness need not be
the unique best bid. This second-price example is a special case.
Assume quasi-linear private values, a report-independent feasible set $\mathcal K$, enforceable transfers without binding budget constraints, and exact maximization of reported total value. Write $\hat v_j(k)=v_j(k;\hat\theta_j)$ for the value induced by report
$\hat\theta_j$, and let $k^*(\hat\theta)$ maximize the sum of reported values.
With a fixed tie rule, the Clarke pivot payment is~\pdcite{vickrey1961counterspeculation,clarke1971multipart,groves1973incentives}
  \begin{equation}
  p_i(\hat\theta)=\max_{k\in\mathcal K}\sum_{j\ne i}\hat v_j(k)
  -\sum_{j\ne i}\hat v_j(k^*(\hat\theta)).
  \end{equation}
  Truthful reporting is \emph{weakly} dominant: the agent's utility is total reported welfare with its own value reported truthfully, minus a term independent of its report. A truthful report selects a welfare maximizer; another report cannot improve it. Equality is possible. In a second-price auction, value 10 and highest rival bid 5 give utility 5 when bidding either 9 or 10.

  \item \textbf{The bilateral-trade boundary:}\enspace Under independent private buyer and seller values with positive densities on intervals whose interiors overlap, Myerson--Satterthwaite rules out simultaneous Bayesian incentive compatibility, interim individual rationality, ex-post efficient trade and expected budget balance without outside subsidy~\pdcite{myerson1983efficient}. Efficiency is essential: no trade and no transfers satisfy the other requirements. VCG does not generally solve budget balance; in bilateral trade it can require a subsidy.

\end{enumerate}

\subsection{Social Choice Theory and Arrow's Impossibility Theorem}

When a swarm of autonomous agents must aggregate individual preference orderings over candidate actions (e.g., selecting an architectural proposal, choosing a refactoring strategy), social choice theory governs the aggregation rules.

\begin{pdexample}{Three consistent voters, one cyclic majority}
The rankings are $A\succ B\succ C$, $B\succ C\succ A$, and $C\succ A\succ B$.
Two voters put $A$ above $B$, two put $B$ above $C$, and two put $C$ above $A$.
Each voter is consistent; the pairwise majority relation cycles. This example
refutes universal transitivity of majority ranking. It does not prove the more
general theorem below, which concerns every rule satisfying its assumptions.
\end{pdexample}

Let $A$ be a finite set of alternatives and $N = \{1, \dots, n\}$ be the set of voting agents. Each agent $i$ has a strict preference ordering $\succ_i \in \mathcal{L}(A)$. A \textbf{Social Welfare Function (SWF)} is a mapping:
\begin{equation}
F: \mathcal{L}(A)^n \;\to\; \mathcal{L}(A)
\end{equation}
that maps a profile of individual rankings $(\succ_1, \dots, \succ_n)$ into a collective ranking $\succ$.

In 1951, Kenneth Arrow proved that a social welfare function on an unrestricted domain cannot satisfy the following criteria together with non-dictatorship~\pdcite{arrow1951social}:

\begin{theorem}[Arrow's Impossibility Theorem (1951)]\label{thm:arrow}
If $|A| \ge 3$, any Social Welfare Function $F: \mathcal{L}(A)^n \to \mathcal{L}(A)$ that satisfies:
\begin{enumerate}[leftmargin=1.5em,itemsep=1.5pt]
  \item \textbf{Unrestricted Domain (U):} $F$ is defined for all logically possible preference profiles.
  \item \textbf{Weak Pareto (P):} If $\forall i \in N, x \succ_i y$, then $x \succ y$.
  \item \textbf{Independence of Irrelevant Alternatives (IIA):} For any two alternatives $x, y$, the collective ranking between $x$ and $y$ depends solely on individual preferences between $x$ and $y$.
\end{enumerate}
is necessarily a \textbf{Dictatorship (D)}: $\exists i \in N$ such that $\forall x, y \in A, \; x \succ_i y \implies x \succ y$.
\end{theorem}

\begin{figure}[tp]
  \centering
  % Reader question: can consistent individual rankings yield a majority cycle?
% Claim: this profile gives a Condorcet cycle; cardinal sums use extra assumptions.
% Evidence: three explicit rankings, compatible utility bids and their totals.
% Distinguish: majority non-transitivity, Arrow hypotheses, cardinal comparability.
% Grammar: paired triangles on the upstream grid, not an unconditional escape.
% Acceptance: recount each 2-to-1 vote and all three cardinal sums.
\noindent\begin{minipage}{11.1cm}\centering
\pdfigurehue{pdcobalt}%
\begin{tikzpicture}[pd figure,x=1cm,y=1cm]
  % Paired triangles have the same geometry and native label size.
  \def\PX{5.65}\def\xA{2.55}\def\xB{1.15}\def\xC{4.00}
  \def\ytop{1.65}\def\ybase{-0.25}
  \node[pd title,anchor=west] at (.10,4.45) {ordinal ranking};
  \node[pd title,anchor=west] at (\PX,4.45) {cardinal utility};
  \foreach \i/\s in {0/{$A\succ B\succ C$},1/{$B\succ C\succ A$},2/{$C\succ A\succ B$}}{
    \node[pd label,anchor=west] at (.10,{3.90-\i*.45}) {agent \number\numexpr\i+1\relax:\ \s};}
  \foreach \i/\s in {0/{$A\,.9$\ $B\,.5$\ $C\,.3$},1/{$A\,.4$\ $B\,.9$\ $C\,.8$},2/{$A\,.6$\ $B\,.2$\ $C\,.9$}}{
    \node[pd label,anchor=west] at (\PX,{3.90-\i*.45}) {agent \number\numexpr\i+1\relax:\ \s};}
  \draw[pd hairline] (\PX,2.80) -- (10.80,2.80);
  \node[pd title,anchor=west] at (\PX,2.50) {$\Sigma$: $A\,1.9$\ $B\,1.6$\ $C\,2.0$};
  \node[pd automaton state] (La) at (\xA,\ytop) {$A$};
  \node[pd automaton state] (Lb) at (\xB,\ybase) {$B$};
  \node[pd automaton state] (Lc) at (\xC,\ybase) {$C$};
  \draw[pd breach arrow] (La) -- (Lb);
  \draw[pd breach arrow] (Lb) -- (Lc);
  \draw[pd breach arrow] (Lc) -- (La);
  \node[pd breach label] at (\xA,-1.15) {$A\succ B\succ C\succ A$: no winner};
  \node[pd automaton state] (Ra) at ({\PX+\xA},\ytop) {$A$};
  \node[pd automaton state] (Rb) at ({\PX+\xB},\ybase) {$B$};
  \node[pd automaton focus] (Rc) at ({\PX+\xC},\ybase) {$C$};
  \draw[pd focus arrow] (Rc) -- (Ra);
  \draw[pd focus arrow] (Ra) -- (Rb);
  \draw[pd focus arrow] (Rc) -- (Rb);
  \node[pd focus label] at ({\PX+\xA},-1.15) {$C\succ A\succ B$: total order};
\end{tikzpicture}
\par\smallskip
{\pdfiglabel\renewcommand{\arraystretch}{1.15}%
\begin{tabularx}{\linewidth}{@{}p{3.25cm}>{\raggedright\arraybackslash}X@{}}
\toprule
\textbf{Condition} & \textbf{Ordinal scope}\\
\midrule
U: unrestricted domain & All strict ranking profiles are allowed.\\
P: weak Pareto & Unanimous pairwise preferences are respected.\\
IIA: independence & A pair depends only on its individual comparisons.\\
D: dictatorship & Forced by U, P, IIA and complete transitivity.\\
\bottomrule
\end{tabularx}}
\par\smallskip
{\pdfiglabelitalic\raggedright Cycles are possible, not inevitable. Arrow assumes a complete,
transitive social order. Cardinal summation uses additional scale and comparability
assumptions; it gives no unconditional governance guarantee.\par}
\end{minipage}
%
  \pdwidecaption{\textbf{A majority-cycle witness.} Three individually consistent rankings yield pairwise majorities $A\succ B$, $B\succ C$, and $C\succ A$. This is one possible profile. Arrow's theorem separately requires unrestricted domain, Pareto and IIA for complete and transitive collective rankings.}[fig:arrows-impossibility]
\end{figure}

\subsubsection{What Aggregation Theorems Require}
\label{sec:why-arrows-matters}

Preference aggregation and safety certification answer different questions. Arrow's theorem~\pdcite{arrow1951social} constrains how complete rankings are combined. It does not prohibit choosing a candidate by voting or prove that every vote deadlocks.
\begin{enumerate}[leftmargin=1.8em,itemsep=2pt]
  \item \textbf{Cycles are possible:} The three rankings $A\succ B\succ C$, $B\succ C\succ A$ and $C\succ A\succ B$ produce the majority cycle $A\succ B\succ C\succ A$. Unanimous rankings do not. A procedure can still return a winner by a declared agenda or tie rule, while giving up other properties.
  \item \textbf{Strategy-proof choice has hypotheses:} Gibbard--Satterthwaite applies to deterministic, onto, single-winner choice rules over at least three alternatives, with unrestricted strict preferences. A rule satisfying these conditions is either dictatorial or manipulable at some profile~\pdcite{gibbard1973manipulation,satterthwaite1975strategy}. It does not say every agent can manipulate every profile.
  \item \textbf{The missing condition matters:} Pareto and a transitive collective ranking alone do not imply dictatorship. A Borda ranking with a fixed tie rule satisfies both on strict preferences but violates IIA. Arrow's impossibility requires IIA as well as unrestricted domain and Pareto.
\end{enumerate}

An agent system must state which aggregation properties its task needs. Restricted preference domains, randomized choice, cardinal reports with transfers, and ordinary voting rules have different guarantees. Cardinal bids and collateral alone do not establish truthful elicitation; VCG needs the assumptions just stated. For mechanically specified admission predicates, a test or proof checks that predicate. Passing a finite test suite does not establish all semantic correctness, and subjective judgments need an explicit decision and appeal rule.

\subsection{The Ostrom Institutional Commons}

When autonomous agents share developer workstation resources (TCP ports, CPU threads, disk space, git branches), they interact within a \textbf{Common-Pool Resource (CPR)}. Elinor Ostrom's Nobel prize-winning institutional economics~\pdcite{ostrom1990governing} disproved the dogma that common resources must inevitably suffer the ``tragedy of the commons'' or be privatized.

Ostrom established eight core design principles for sustainable, self-governing institutions:
\begin{enumerate}[leftmargin=1.8em,itemsep=2pt]
  \item \textbf{Clearly Defined Boundaries:} Clear demarcation of who has access to which resources.
  \item \textbf{Congruence Between Appropriation Rules and Local Conditions:} Resource consumption rules match environmental scarcity.
  \item \textbf{Collective-Choice Arrangements:} Individuals affected by operational rules can participate in modifying them.
  \item \textbf{Active Monitoring of Resource Condition and User Behavior:} Monitors are accountable to users.
  \item \textbf{Graduated Sanctions:} Violators face escalating penalties rather than binary expulsion.
  \item \textbf{Low-Cost Conflict Resolution Mechanisms:} Rapid, local arbitration of disputes.
  \item \textbf{Minimal Recognition of Rights to Organize:} External authorities do not override self-governance.
  \item \textbf{Nested Enterprises:} Governance activities are organized in multiple layers of nested institutions.
\end{enumerate}

Chapter~\ref{chap:bonded} uses these principles as design questions for the \textbf{Bonded Commons}: who may use a resource, who observes violations, and how can a ruling be challenged? Its conditional economic model does not establish that an implemented institution will be self-governing or resilient.

\section{The Physics of Modern LLM Execution: What an Operating Kernel Must Know}\label{sec:prereq-llm-physics}

To govern large language models, an operating system must understand the physical and computational realities of their execution. This section establishes the fundamental physics of modern LLMs: attention mechanics, memory bandwidth bottlenecks, training alignments, stochastic entropy, and context economics.

\subsection{Transformer Mechanics: Self-Attention, Quadratic Complexity, and the KV-Cache}

\begin{pdexample}{What an attention weight does}
For one query, suppose two admissible positions have equal scores and scalar
values 2 and 4. Softmax converts scores $z_j$ to weights
$e^{z_j}/\sum_k e^{z_k}$; equal scores therefore give weights $1/2,1/2$.
The weighted output is $\tfrac12(2)+\tfrac12(4)=3$. This arithmetic is a mixing
operation, not evidence that either input statement is true. The cost of
forming many such scores is the point of the following model.
\end{pdexample}

The Transformer~\pdcite{vaswani2017attention} is a widely used architecture. Let $X \in \mathbb{R}^{L \times d}$ be an input sequence of $L$ tokens projected into a $d$-dimensional embedding space. For one attention head, the input is mapped via learned weight matrices $W_Q, W_K, W_V \in \mathbb{R}^{d \times d_k}$ into Queries ($Q$), Keys ($K$), and Values ($V$):
\begin{equation}
Q = X W_Q, \quad K = X W_K, \quad V = X W_V
\end{equation}
Softmax is applied separately to each score row. A causal mask excludes future positions in autoregressive use. Multiple heads perform separate projections and combine their outputs; the following expression describes one head:
\begin{equation}\label{eq:attention-formula}
\mathrm{Attention}(Q, K, V) \;=\; \mathrm{softmax}\left( \frac{Q K^T}{\sqrt{d_k}} \right) V
\end{equation}

\subsubsection{The Quadratic Attention Wall}
The matrix product $Q K^T$ forms an $L \times L$ attention matrix representing the pairwise affinity between every token in the sequence. For dense attention that materializes the score matrix, the asymptotic costs are:
\begin{equation}
\mathrm{Compute} \;=\; O(L^2 \cdot d), \qquad \mathrm{Memory} \;=\; O(L^2 + L \cdot d)
\end{equation}
While FlashAttention~\pdcite{dao2022flashattention} uses GPU SRAM tiling to avoid materializing the full $L \times L$ matrix in High-Bandwidth Memory (HBM), the underlying FLOP count remains quadratic in context length $L$. For an agent maintaining an uncompacted context of $L = 128{,}000$ tokens, a single attention pass requires evaluating $\approx 1.63 \times 10^{10}$ token-pair affinities per attention head.

\subsubsection{The Key-Value (KV) Cache and Memory Bandwidth Bottlenecks}
In autoregressive text generation, consider a dense decoder generating one token at a time: $x_{t} \sim P_\theta(x_t \mid x_{<t})$. Because previous tokens do not change, recomputing their Key and Value representations at every decoding step would result in $O(L^3)$ aggregate compute. Production inference engines (such as vLLM~\pdcite{kwon2023vllm}) therefore cache Key and Value vectors in high-bandwidth GPU memory (HBM).

For a transformer with $N_l$ layers, $N_h$ key-value heads, dimension $d_h$, sequence length $L$, and precision $b_p$ bytes (typically 2 for FP16/BF16), the KV-cache memory consumption per active agent sequence is:
\begin{equation}\label{eq:kv-cache-full}
\mathrm{Memory}_{\mathrm{KV}} \;=\; 2 \times N_l \times N_h \times d_h \times L \times b_p \quad \text{bytes}
\end{equation}

Consider an illustrative 70-billion parameter model configuration ($N_l = 80$, $N_h = 8$ using Grouped-Query Attention, $d_h = 128$, $b_p = 2$). At context length $L = 128{,}000$:
\begin{align}
\mathrm{Memory}_{\mathrm{KV}} &\;=\; 2 \times 80 \times 8 \times 128 \times 128{,}000 \times 2 \;\text{bytes} \notag \\
&\;\approx\; 41.94\text{ GB per sequence}
\end{align}
A fleet of 8 concurrent agent workers requires $\approx 335.5\text{ GB}$ of high-bandwidth memory solely to store the KV-cache, exceeding a nominal 320~GB memory budget before allocating a single byte for model weights.

\subsubsection{Arithmetic Intensity: Compute-Bound Prefill vs. Memory-Bound Decoding}
A critical physical reality of GPU execution is \textbf{Arithmetic Intensity}: the ratio of floating-point operations (FLOPs) to memory access (bytes transferred):
\begin{equation}
I \;=\; \frac{\text{FLOPs}}{\text{Bytes Transferred}}
\end{equation}
Prefill and decoding often have different bottlenecks:
\begin{enumerate}[leftmargin=1.8em,itemsep=2pt]
  \item \textbf{Prefill:} Processing many prompt tokens uses matrix multiplications and can achieve high arithmetic intensity. Whether it is compute-bound depends on batch size, sequence length, hardware and the attention implementation.
  \item \textbf{Decoding:} Single-sequence dense decoding often spends substantial bandwidth reading weights and the KV cache. A dense 70-billion-parameter model stored at two bytes per parameter has 140~GB of weights before sharding or compression. Batching, quantization, sparse architectures and speculative decoding change the balance.
\end{enumerate}
Long contexts therefore have a measurable memory and compute cost. The cache formula provides a capacity estimate for the stated configuration; latency and throughput require measurements of the actual inference system.

\subsection{Training Regimes and the Alignment Illusion}

Training changes how proposals are generated. Pretraining learns predictive
patterns; supervised fine-tuning uses demonstrations of desired responses;
preference training uses comparisons or rewards to favor selected behavior
\pdcite{ouyang2022instruct,rafailov2023dpo}. These objectives are useful context,
but their loss functions do not enter any proof in this book.

For Alice's job, training might make a worker more likely to request a test
before claiming success. An admission rule can separately require a receipt
for the specified revision. Improving the request probability and checking the
receipt solve different problems. The distinction is the one introduced by the
counter example: a good proposal is not yet a safely executed transition.


\subsubsection{Alignment Objectives and Execution Invariants}
Preference training can improve observed behavior. Its objective alone is not a proof that every permitted action preserves a host invariant. An apology or instruction remains evidence about the policy's behavior; an admission rule constrains which effects can occur. The distinction depends on the action interface and the verified controller, not on whether the proposer is neural or symbolic.

\subsection{Sampling and Conditional Failure Hazards}

For finite logits $z$ and temperature $T>0$, softmax assigns
\begin{equation}
P(w_i\mid x_{<t})=\frac{\exp(z_i/T)}{\sum_j\exp(z_j/T)}.
\end{equation}
Greedy decoding at $T=0$ is a separate argmax rule with a specified tie policy. For nucleus sampling, order tokens so $P(w_1)\ge\cdots\ge P(w_m)$, fix a tie order, and let $0<p\le1$. Following \href{https://arxiv.org/abs/1904.09751}{Holtzman et al.'s nucleus rule}, retain the shortest prefix reaching probability $p$ and renormalize:
\begin{equation}
k_p=\min\left\{k:\sum_{j=1}^{k}P(w_j)\ge p\right\},\qquad
V^{(p)}=\{w_1,\ldots,w_{k_p}\}.
\end{equation}
With probabilities $(0.95,0.05)$ and $p=0.9$, the nucleus is $\{w_1\}$. Truncation, action masks and tool validation can exclude unsafe actions even when token generation is stochastic.

\begin{pdexample}{Two steps of conditional risk}
Suppose a run survives its first step with probability 0.9 and, conditional on
that survival, its second with probability 0.8. It survives both with probability
$0.9\times0.8=0.72$. Independence is unnecessary because the second probability
is already conditional. The theorem uses the same multiplication with upper
bounds on each conditional survival probability.
\end{pdexample}

\begin{theorem}[Survival under conditional failure hazards]\label{thm:non-determinism}
Let $S_t$ mean that execution has remained safe through step $t$, with $S_0$ certain. Assume that after every reachable safe history before step $t$, the probability that the \emph{admitted effect} violates the invariant is at least $\epsilon_t\in[0,1]$. Then
\begin{equation}
\Pr(S_N)\le\prod_{t=1}^{N}(1-\epsilon_t).
\end{equation}
If one constant $\epsilon>0$ lower-bounds every $\epsilon_t$, then $\Pr(S_N)\le(1-\epsilon)^N\to0$. No independence assumption is required.
\end{theorem}
\noindent\textit{Proof idea.} Bound survival at the next step conditional on having survived so far. Averaging the history-wise bound gives $\Pr(S_t\mid S_{t-1})\le1-\epsilon_t$. Multiply these conditional probabilities. If any $S_t$ already has probability zero, the conclusion is immediate.\hfill$\square$

The hazard assumption must be measured or otherwise justified; positive temperature does not imply it. A policy restricted to safe actions has zero failure probability. Even positive hazards at every step need not force eventual failure: independent hazards $\epsilon_t=2^{-t-2}$ have total at most $1/4$, so the union bound leaves survival probability at least $3/4$. This theorem explains how persistent admitted risk accumulates. It neither proves that every model must fail nor equates a proposed unsafe token with an executed unsafe effect.

\subsection{The Context Window as Volatile, Lossy RAM}

A prompt can contain an exact string without ensuring that a generated answer
will retrieve or use it correctly. This distinguishes retaining bytes from
successfully using their contents.

\subsubsection{Position and retrieval}
Liu et al.~\pdcite{liu2024lost} measured sensitivity to the position of relevant
information in the models and tasks they tested. Performance was often better
at the beginning or end than in the middle. These are task-performance results,
not a universal percentage loss or a direct measurement of attention weights.
The engineering question is whether the actual system retrieves the evidence
needed for the current decision.

\subsubsection{Compaction and retained evidence}
\begin{pdexample}{Two histories, one summary}
History A says ``test failed on revision 7''; history B says ``test not run on
revision 7.'' A summary that represents both as ``testing pending'' cannot
recover which history occurred: two distinguishable inputs have one output. This is loss by construction, even if the
summary itself contains no false sentence. Retaining a receipt and linking it
to the summary permits a later reader to distinguish them. Chapter 4 asks how
to allocate scarce summary space between continuing work and checking it.
\end{pdexample}
Compaction can preserve particular facts exactly; loss is not inevitable for
every query. Its contract must say which distinctions survive and which must
be recovered from retained evidence. The task's authoritative obligations and
resource permissions cannot depend solely on a model remembering them.


\begin{pdclaim}{Design invariant}{Externalized Durable State}\label{req:external-state}
An agent operating kernel must never rely on the model's context window as the system of record. Authoritative locks, capabilities, claims, balances and evidence references need durable storage with a stated consistency and recovery contract. Chapter 1 studies a transactional database implementation; other implementations must discharge the same obligations. Prompt buffers contain working views of that state.
\end{pdclaim}

\section{Modern Coding Agent Architectures, Planning, and Search}\label{sec:prereq-coding-agents}

We now examine how modern LLM-based coding agents interface with host environments, structure deliberation, and execute search over complex problem spaces.

\subsection{The Contemporary Landscape: Systems Comparison}

Coding harnesses combine proposal generation with host tools. Compare an
execution path by what it can establish, rather than treating a product name
as a fixed security boundary.
\begin{table}[H]
\centering\small
\begin{tabularx}{\textwidth}{@{}lX@{}}
\toprule
\textbf{Boundary} & \textbf{Question to answer for the actual path}\\
\midrule
Isolation & Which files, network destinations and subprocesses are reachable?\\
State & Which facts survive compaction and process replacement?\\
Admission & Can any actuator bypass permission or version checks?\\
Recovery & Can a retry duplicate a committed effect?\\
Evidence & Which artifact and revision does a success report identify?\\
\bottomrule
\end{tabularx}
\caption{Questions for inspecting a coding harness. Later chapters supply the contracts.}
\label{tab:coding-agents}
\end{table}

\subsection{Closed-Loop Deliberation and Partial Observation}
\paragraph{Example: A timeout is not a failed write.}
A worker sends a write request and receives no response. In one possible state,
the request never arrived. In another, it committed and its response was lost.
Retrying blindly is safe in neither generality nor by virtue of persistence:
a payment might be duplicated. The worker needs an operation identifier and
an authoritative lookup, or an explicitly idempotent operation (repeating it
has the same intended effect). Chapter 1 supplies the corresponding recovery
obligations. This is why the earlier failed-plan example distinguished a failed
observation from the state of the world.


A partially observable Markov decision process (POMDP) is one mathematical
model for this distinction. Its tuple
$\langle\mathcal S,\mathcal A,\mathcal T,\mathcal R,\Omega,\mathcal O,\gamma\rangle$
names states, actions, transition probabilities, rewards, observations,
observation probabilities and a discount factor $0\le\gamma<1$.
A belief state $b_t(s)=\Pr(s_t=s\mid h_t)$ is a probability distribution over
states conditional on the action--observation history $h_t$. The timeout
leaves probability on both states until evidence distinguishes them.

Calling a harness partially observable does not show that it computes a
Bayesian belief state or solves a POMDP optimally. A closed loop means that
observations can change subsequent actions. It permits correction; it does
not guarantee recovery or make external effects reversible.


\subsection{Deliberation, Planning, and Search Frameworks}

Modern agents employ structured search and planning algorithms across candidate execution paths~\pdcite{wang2023agentsurvey}:

\begin{enumerate}[leftmargin=1.8em,itemsep=3pt]
  \item \textbf{ReAct (Reasoning and Acting)~\pdcite{yao2022react}:}\enspace Interleaves natural language reasoning traces (``Thoughts'') with structured tool executions (``Actions'') and environment feedback (``Observations''). An action can gather evidence that changes the next reasoning step. This interleaving permits inquiry; it does not guarantee that every uncertainty triggers a tool call.
  
  \item \textbf{Tree of Thoughts (ToT)~\pdcite{yao2023tree}:}\enspace Explores alternative intermediate reasoning steps in a tree. The studied searches proceed breadth-first (comparing a frontier of candidates) or depth-first (following a candidate and backtracking). Model evaluations guide exploration; they are heuristics, not proofs that a branch is correct.
  
  \item \textbf{Reflexion~\pdcite{shinn2023reflexion}:}\enspace Uses linguistic feedback without updating the model weights. The agent reflects on task feedback, retains that reflection in an episodic memory buffer, and consults it in later attempts. For example, a failed test can prompt a reminder to check the precondition that was missed.

  \item \textbf{Voyager Lifelong Skill Learning~\pdcite{wang2023voyager}:}\enspace In Minecraft, an agent accumulates an \textbf{executable skill library}: code for behaviors that later tasks can retrieve and compose. Environment feedback, execution errors and self-verification guide program improvement. The example shows how procedural memory can outlast a prompt; using the pattern for software work still requires checking dependencies, permitted effects and task-specific correctness.

  \item \textbf{TDAG~\pdcite{wang2024tdag}:}\enspace Task Decomposition and Agent Generation dynamically decomposes a task and generates specialized subagents. Wang et al. evaluate it on ItineraryBench, a travel-planning benchmark. This supports an adaptive planning pattern; it is not a theorem about eliminating dependency cycles or serializing software mutations.

  \item \textbf{HyperTree Planning~\pdcite{gui2025hypertree}:}\enspace Gui et al. construct and iteratively refine hierarchical planning outlines, evaluated on TravelPlanner. The structure supports decomposition of constrained plans. Its reported benchmark gains do not establish that software agents avoid local minima or produce correct repository changes.
\end{enumerate}

\subsection{Cognitive Control, Test-Time Compute, and Process Reward Models}

Inference search can explore alternative reasoning chains scored by \textbf{Process Reward Models (PRMs)}. For input $x$, let $\tau=(c_1,\ldots,c_T)$ be a candidate chain and let the PRM assign each step a score in $[0,1]$ given its prefix. One possible ranking heuristic multiplies these scores:
\begin{equation}
\mathrm{Score}(\tau) \;=\; \prod_{t=1}^T \mathrm{PRM}(c_t \mid c_{<t}, x)
\end{equation}
This product is a joint correctness probability only when its factors are the appropriate calibrated conditional probabilities. For instance, scores 0.9 and 0.8 multiply to 0.72; without that calibration, the number does not mean a 72\% chance of a correct answer. Otherwise it is a heuristic score. A search policy may backtrack when a score falls below a chosen threshold.

\subsection{Multi-Agent Formal Verification and Benchmarking}

Evaluating multi-agent software engineering requires moving beyond synthetic static benchmarks (such as HumanEval) toward interactive, sandboxed evaluation suites:
\begin{itemize}[leftmargin=1.8em,itemsep=2pt]
  \item \textbf{AgentBench~\pdcite{liu2023agentbench}:} A multi-dimensional benchmark evaluating LLMs as autonomous agents across operating system bash environments, databases, web browsing, and multi-agent coordination games.
  \item \textbf{HELM (Holistic Evaluation of Language Models)~\pdcite{liang2022helm}:} Standardized multi-metric evaluation across accuracy, calibration, robustness, fairness, and toxicity.
  \item \textbf{FormalJudge~\pdcite{zhou2026formaljudge}:} Zhou et al. translate natural-language requirements into atomic constraints and use Dafny specifications with Z3 to assess compliance. Their agent-oversight benchmarks report improved detection. A proof concerns the encoded specification; correctness of its translation and of its connection to observed behavior remains a separate obligation.
\end{itemize}

\section{Tools, Protocols, Skills, and Hooks}\label{sec:prereq-hypervisors}

To interact safely with the host environment, an agent harness must govern tool dispatch, encapsulate domain knowledge, and enforce runtime invariants.

\subsection{Model Context Protocol (MCP) as the Substrate}

The \href{https://modelcontextprotocol.io/specification/2025-11-25/server/tools}{Model Context Protocol (MCP), specification 2025-11-25}, standardizes client--server interfaces for tools and context. It does not by itself confine effects:
\begin{enumerate}[leftmargin=1.8em,itemsep=2pt]
  \item \textbf{Tools:} Server-exposed procedures invoked through the harness ($\mathtt{tools/call}$), with JSON-Schema inputs and results that can contain content blocks or structured data.
  \item \textbf{Resources:} Passive data endpoints ($\mathtt{resources/read}$) exposing context such as local file contents, system logs, or git diffs.
  \item \textbf{Prompts:} Reusable prompt templates parameterizing multi-step workflows.
\end{enumerate}

\subsection{Context Economics and Progressive Skill Disclosure}

Progressive disclosure keeps an index in context and loads detailed instructions when needed. Its benefit depends on actual retrieval and task frequency:
\begin{enumerate}[leftmargin=1.8em,itemsep=3pt]
  \item \textbf{Index:} A concise name and description identify each skill. For illustration, 500 descriptions averaging 30 tokens consume 15,000 tokens, or $7.5\%$ of a 200,000-token window. This is a sizing example, not a measured activation rate.
  \item \textbf{Instructions:} A selected skill supplies operational guidance, parameters and recovery rules. Its context cost is the number of instruction tokens actually loaded.
  \item \textbf{References and execution:} The harness reads relevant reference excerpts or runs a script without loading all its source into the model. Tool arguments, outputs and errors still consume context. Executing a script is not a zero-token interaction.
\end{enumerate}
If the index costs $b_0$ tokens and asset $i$ costs $b_i$ when loaded with probability $p_i$, one loading round has expected cost $b_0+\sum_i p_i b_i$. This follows by linearity of expectation without requiring independence. It gives no exponential saving unless a measured access distribution justifies one.

\begin{figure}[tp]
  \centering
  % Reader question: Which material must enter the model context now?
% Claim: Staged loading separates discovery, instructions and execution.
% Evidence: Four artifact roles, with conditional retrieval.
% Distinguish: Discovery metadata, task text, references and host code.
% Grammar: A retrieval dependency graph, not invented frequency measurements.
% Acceptance: Identify which text is loaded only when needed.
\pdfigurehue{pdcobalt}%
\begin{tikzpicture}[pd figure,x=1cm,y=1cm]

  \node[pd artifact,text width=6.0cm,align=center] (meta) at (5.2,0) {Discovery metadata\\Name, purpose and activation conditions};
  \node[pd focus state,text width=6.0cm,align=center] (skill) at (5.2,-1.9) {Selected skill instructions\\Load when the task needs them};
  \node[pd artifact,text width=3.7cm,align=center] (refs) at (2.4,-3.9) {Referenced material\\Retrieve relevant sections};
  \node[pd state,text width=3.7cm,align=center] (code) at (8.0,-3.9) {Host-executed scripts\\Return bounded results};
  \draw[pd arrow] (meta) -- (skill);
  \draw[pd arrow] (skill) -- (refs);
  \draw[pd arrow] (skill) -- (code);
  \node[pd note,text width=9.4cm,align=left] at (5.2,-5.5) {Context savings depend on what is retrieved and what the tools return. This architecture alone does not establish a usage frequency or an exponential reduction.};

\end{tikzpicture}
%
  \pdwidecaption{\textbf{Staged loading.} Keep an index, load selected instructions, then retrieve references or execute scripts as needed. Savings depend on actual access patterns and returned output; the tiers imply no universal activation frequencies.}[fig:progressive-skills-iceberg]
\end{figure}

\vspace{10pt}
\noindent\begin{minipage}{\linewidth}
  \centering
  \pdswissfigplate{fig-ch0-skills-logarithmic}{0.88\linewidth}%
  \par\vspace{8pt}%
  {\color{pdcobalt}\rule{\linewidth}{1.2pt}}\par\vspace{5pt}%
  \raggedright
  {\color{pdswissred}\rule{4.5pt}{4.5pt}}\enspace%
  {\pdgrotesk\bfseries\scshape\fontsize{9}{11}\selectfont\color{hhink}%
    Architectural Plate 0.D\enspace\textperiodcentered\enspace Progressive Skill Disclosure Economics}\par\vspace{4pt}%
  {\pdgrotesk\itshape\fontsize{9.5}{13}\selectfont\color{hhink}%
    ``A compact index directs selective loading; script outputs and reference excerpts determine the context actually consumed.''}\par\vspace{6pt}%
  {\color{hhgray!30}\rule{\linewidth}{0.4pt}}
\end{minipage}
\vspace{10pt}

\subsection{Lifecycle Hooks and Supervisory Control}

\paragraph{Example: A write can be refused; a crash must be accommodated.}
Suppose the next event is either a requested write or a power failure. A gate
can refuse the write. It cannot refuse the power failure. If safety requires
recoverable obligations after either event, the current state must already
retain enough information to survive the failure. Choosing ``deny'' at a
post-crash hook is too late. This is the recovery half of the opening problem.


Ramadge--Wonham supervisory control~\pdcite{ramadge1987supervisory} models
which events a supervisor can disable. A lifecycle hook can implement part
of such a supervisor only when the modeled control and observation assumptions
hold for its actual effect paths.

Let unconstrained agent execution be modeled as a discrete-event generator $G = \langle X, \Sigma, \delta, x_0, X_m \rangle$, where $X$ is the state set, $x_0$ the initial state, $\delta$ the partial state-transition function, $X_m$ the marked completion states, and events $\Sigma = \Sigma_c \cup \Sigma_u$ are partitioned into controllable events $\Sigma_c$ (tool invocations, file writes, network requests) and uncontrollable events $\Sigma_u$ (model token generation, host hardware interrupts, external process termination). A trace is a finite event sequence. $L(G)$ contains the possible traces; $\mathcal P(\Sigma)$ denotes event subsets. A supervisor $S:L(G)\to\mathcal P(\Sigma)$ selects enabled events after each trace under full observation, always allowing uncontrollable events. If $K$ contains the safe traces and is prefix-closed (every initial segment of a safe trace is safe), the desired controlled language satisfies:
\begin{equation}
L(S / G) \;\subseteq\; K \;\subseteq\; L(G)
\end{equation}
where $K$ is a prefix-closed language of safe traces. This inclusion is a design requirement, not an automatic property of a hook. For full observation, enforcing exactly $K$ requires controllability: $K\Sigma_u\cap L(G)\subseteq K$. Otherwise one must find an enforceable sublanguage, or report that the initial state has no safe continuation. Partial observation adds observability requirements; completion adds a separate nonblocking obligation.

Figure~\ref{fig:lifecycle-hooks-engines} illustrates pre-effect admission and
post-effect observation. The event names follow Claude Code's hook
reference~\pdcite{anthropic2026hooks}; the drawing is a teaching model, not a
complete implementation map. Two useful levels are:
\begin{enumerate}[leftmargin=1.8em,itemsep=2pt]
  \item \textbf{The Outer Session \& Turn Loop:} Encompasses initial context injection (\texttt{SessionStart}), prompt interception (\texttt{UserPromptSubmit}), context pressure management (\texttt{PreCompact}), turn completion checks (\texttt{Stop}), error notification (\texttt{StopFailure}), and session teardown (\texttt{SessionEnd}).
  \item \textbf{The Inner Agentic Tool Loop:} Governs the iterative reasoning-action cycle: pre-tool interception and deterministic vetoes (\texttt{PreToolUse}), sandboxed execution, and reactive post-tool sanitization (\texttt{PostToolUse}).
\end{enumerate}

\begin{figure}[tp]
  \centering
  % Reader question: which hook precedes effects and which inspects results?
% Claim: a pre-tool veto can prevent a mediated call; post-tool checks run after it.
% Evidence: upstream nested session/turn and agentic tool loops.
% Distinguish: admission, confined effects, result inspection and closeout.
% Grammar: retain the nested loop and named interception points.
% Acceptance: distinguish a pre-effect veto from limited after-effect compensation.
\noindent\begin{minipage}{11.1cm}\centering
\pdfigurehue{pdcobalt}%
\begin{tikzpicture}[pd figure,x=1cm,y=1cm]
  % The outer spine and inner tool cycle keep independent, bounded lanes.
  \def\sx{2.70}\def\ix{8.50}
  \node[pd focus state,minimum width=2.5cm] (O1) at (\sx,4.20) {prompt received};
  \node[pd focus state,minimum width=2.5cm] (O2) at (\sx,2.80) {turn executing};
  \node[pd focus state,minimum width=2.5cm] (O3) at (\sx,1.40) {turn reviewed};
  \node[pd terminal,minimum width=2.5cm] (O4) at (\sx,0) {session closed};
  \draw[pd initial] (\sx,5.45) -- (O1);
  \node[pd label] at (\sx,5.70) {SessionStart};
  \draw[pd focus arrow] (O1) -- (O2);
  \node[pd badge] at (2.70,3.50) {1};
  \draw[pd focus arrow] (O3) -- (O4);
  \node[pd label,anchor=west] at (2.95,.70) {SessionEnd};
  \draw[pd focus arrow,rounded corners=3pt] (O3.west) -- (.25,1.40) -- (.25,2.80) -- (O2.west);
  \node[pd focus label,rotate=90] at (.70,2.10) {next turn};
  \node[pd artifact,text width=1.5cm,align=center] (PC) at (5.00,4.25) {Pre\\Compact};
  \draw[pd thin arrow,<->] (O2.north east) -- (PC.south west);
  \node[pd badge] at (4.60,3.40) {2};
  \draw[pd trust boundary] (5.95,1.75) rectangle (10.85,4.85);
  \node[pd kind] at (8.40,4.55) {inner agentic tool loop};
  \node[pd focus state,minimum width=2.2cm] (I1) at (\ix,2.40) {deliberate};
  \node[pd focus state,minimum width=2.2cm] (I2) at (\ix,3.80) {sandboxed};
  \draw[pd focus arrow] (I1.west) -- (7.00,2.40) -- (7.00,3.80) -- (I2.west);
  \draw[pd focus arrow] (I2.east) -- (10.10,3.80) -- (10.10,2.40) -- (I1.east);
  \node[pd badge] at (7.00,3.10) {3};
  \node[pd badge] at (10.10,3.10) {4};
  \draw[pd thin arrow] (O2.east) -- (5.25,2.80) -- (5.25,2.40) -- (I1.west);
  \draw[pd focus arrow,rounded corners=3pt] (I1.south) -- (8.50,1.15) -- (4.40,1.15) -- (O3.east);
  \node[pd badge] at (6.15,1.15) {5};
  \node[pd breach state,minimum width=1.8cm] (BX) at (8.50,-.45) {blocked};
  \draw[pd breach arrow] (I1.east) -- (10.80,2.40) -- (10.80,-.45) -- (BX.east);
  \node[pd breach label] at (8.50,-1.05) {exit 2};
\end{tikzpicture}
\par\smallskip
{\pdfiglabel\renewcommand{\arraystretch}{1.15}%
\begin{tabularx}{\linewidth}{@{}c p{2.8cm}p{1.05cm}>{\raggedright\arraybackslash}X@{}}
\toprule
 & \textbf{Hook} & \textbf{Loop} & \textbf{Fires and may}\\
\midrule
1 & \texttt{UserPromptSubmit} & outer & Inject or reject context before deliberation.\\
2 & \texttt{PreCompact} & outer & Record a recovery checkpoint before compaction.\\
3 & \texttt{PreToolUse} & inner & Veto a call before its mediated effect.\\
4 & \texttt{PostToolUse} & inner & Record results after effects; compensation may be limited.\\
5 & \texttt{Stop} & outer & Check closeout; request more work if needed.\\
\bottomrule
\end{tabularx}}
\par\smallskip
{\pdfiglabelitalic\raggedright PreToolUse can veto a mediated call before it runs.
PostToolUse cannot undo irreversible effects. Hook names and veto behavior
depend on the configured host contract.\par}
\end{minipage}
%
  \pdwidecaption{\textbf{Hook ordering.} A pre-execution gate can refuse a controllable action. Post-execution checks inspect effects and evidence. Closeout records the result. Prevention requires every relevant effect path to cross the gate before it occurs.}[fig:lifecycle-hooks-engines]
\end{figure}

\subsubsection{Illustrative Hook Designs}
\label{sec:lifecycle-hook-cases}
The following constructed scenarios separate prevention, recovery, and retained
context. They specify design obligations; they are not measured production runs.

\begin{description}[leftmargin=1.5em,style=sameline,font=\normalfont]
\item[\textbf{Case A: Pre-tool refusal.}]
An agent proposes a destructive command:
\begin{lstlisting}[language=bash,basicstyle=\ttfamily\small,breaklines=true]
git checkout main &&
  git reset --hard origin/main
\end{lstlisting}
A configured admission hook identifies a forbidden protected-branch mutation
and refuses the tool call. In Claude Code, a blocking \texttt{PreToolUse}
command hook can return exit code~2 with a reason on standard error.
Other exit codes depend on the configured failure policy and are not this
explicit veto contract~\pdcite{anthropic2026hooks}. A short
reason is sufficient: ``Protected branch mutation refused.'' A shell-pattern
check by itself does not establish complete mediation: alternate executables,
indirect effects, and other write channels must also be controlled.

\item[\textbf{Case B: Post-tool recovery.}]
A build writes unexpected generated files. A post-tool observer can compare
the resulting paths with a recorded baseline, report the deviation, and invoke
a separately authorized recovery procedure for recoverable files. This happens
after the write; it cannot undo a disclosure or arbitrary remote effect.
Untracked files are not automatically disposable: ownership and recovery
scope must be established before cleanup.

\item[\textbf{Case C: Compaction preparation.}]
Before context compaction, a configured handler records current obligations,
claim identifiers, failing checks, and source locators in a durable checkpoint.
The successor loads and validates that record. The contract must define which
fields are mandatory and how freshness is checked; neither a summary nor a
hook name guarantees their preservation. Secrets remain in the credential
store, with references in the checkpoint. No turn-count or context-percentage
threshold is assumed here.
\end{description}

\vspace{10pt}
\noindent\begin{minipage}{\linewidth}
  \centering
  \pdswissfigplate{fig-ch0-lifecycle-hooks}{0.88\linewidth}%
  \par\vspace{8pt}%
  {\color{pdcobalt}\rule{\linewidth}{1.2pt}}\par\vspace{5pt}%
  \raggedright
  {\color{pdswissred}\rule{4.5pt}{4.5pt}}\enspace%
  {\pdgrotesk\bfseries\scshape\fontsize{9}{11}\selectfont\color{hhink}%
    Architectural Plate 0.E\enspace\textperiodcentered\enspace Dual-Loop Supervisory Interception}\par\vspace{4pt}%
  {\pdgrotesk\itshape\fontsize{9.5}{13}\selectfont\color{hhink}%
    ``Deterministic lifecycle hooks intercept unconstrained generation at binary boundaries, replacing conversational advice with discrete supervisory vetoes.''}\par\vspace{6pt}%
  {\color{hhgray!30}\rule{\linewidth}{0.4pt}}
\end{minipage}
\vspace{10pt}

\section{Multi-Agent Swarm Topologies, Coordination, and Multi-Agent Economics}\label{sec:prereq-swarms}

When scaling from a single agent to a collaborative collective, coordination must transition from direct peer-to-peer communication to structured environmental and institutional governance.

\subsection{Stigmergic Communication: Environmental Coordination}

Communication cost depends on who publishes, who consumes each update, and what state they need. In 1959, French biologist Pierre-Paul Grass{\'e} coined the term \textbf{Stigmergy} (from the Greek \emph{stigma}, mark, and \emph{ergon}, work) to explain how social insects construct intricate, cathedral-like termite mounds without central leadership or point-to-point communication.

\pdmarginfigure{termites}{Cathedral termite mound (\emph{Nasutitermes triodiae}) in Litchfield National Park, Australia. Social insects construct massive architectural edifices through stigmergic environmental modification without central coordination or peer-to-peer communication.}%
In a stigmergic coordination pattern, \textbf{the environment is a communication medium}:
\begin{enumerate}[leftmargin=1.8em,itemsep=2pt]
  \item \textbf{Passive Environmental Modification:} Agents coordinate by leaving persistent, machine-readable traces directly in shared host artifacts: file locks, AST reservation tags, pheromone heat maps, and append-only note streams.
  \item \textbf{Decoupled Sensing:} When Agent~$A$ claims an AST symbol or edits a function, it writes a tuple into the shared kernel tuple space. When Agent~$B$ inspects the repository, it senses Agent~$A$'s trace and autonomously diverts its own planning policy away from the contested file.
\end{enumerate}

If $N$ agents each send one update to all peers, there are $N(N-1)$ deliveries. Publishing each update once to shared storage uses $N$ writes, but if every agent then reads every update, the reads still number $N(N-1)$. Total traffic is linear only under a bounded-read or suitably aggregated subscription policy. Shared traces decouple execution in time; they do not remove the cost of distributing information.

\subsection{Conway's Law in Multi-Agent Swarms and Role-Based SOPs}

In 1968, computer scientist Melvin Conway formulated the foundational law of software architecture~\pdcite{conway1968}:
\begin{quote}
``Organizations which design systems are constrained to produce designs which are copies of the communication structures of these organizations.''
\end{quote}
Conway's observation motivates a design hypothesis for agent swarms; it is not a universal mathematical theorem:
\begin{itemize}[leftmargin=1.8em,itemsep=2pt]
  \item \textbf{The coupling hypothesis:} Unstructured all-to-all communication can make responsibility and interfaces harder to track. It does not logically force a monolithic software design. Evaluate the resulting dependency structure and maintenance outcomes rather than inferring them from the communication graph alone.
  \item \textbf{MetaGPT Role-Based Standard Operating Procedures (SOPs)~\pdcite{hong2024metagpt}:}\enspace Hong et al. study a design that encodes software engineering procedures into roles: Product Managers generate PRDs, Architects synthesize system design diagrams, Engineers write modular code, and QA Reviewers run test suites. Structured documents carry intermediate results. This motivates an interface design; it does not prove that role separation prevents every error.
\end{itemize}

\subsection{Six Sigma Defect Minimization in Agent Swarms}

Patel et al.~\pdcite{patel2026sixsigma} propose task decomposition, redundant micro-agent sampling, and consensus with dynamic sample counts. Their reliability calculation relies on independence of the sampled errors. For odd $n$, independent votes with common error probability $q<1/2$ yield the binary-majority error
\begin{equation}
P_{\mathrm{maj}}=\sum_{j=(n+1)/2}^{n}{n\choose j}q^j(1-q)^{n-j}.
\end{equation}
For $n=5$ and $q=0.05$, this is $0.001158125$, about $0.116\%$. Perfectly correlated errors instead leave error probability $q$, regardless of $n$. These calculations do not certify arbitrary semantic clustering of open-ended answers.

For a sequential pipeline, define $d_k$ as the conditional probability that stage $k$ fails given that all preceding stages passed. Then the chain rule gives
\begin{equation}
\Pr(\text{all stages pass})=\prod_{k=1}^{K}(1-d_k).
\end{equation}
Using unconditional stage error rates in that product would require an independence assumption. Deterministic checks are useful when an admission predicate is available, but this is a separate design choice: a test pass establishes only the tested condition, and consensus alone does not supply a formal oracle.

\subsection{Runtime Task Overlap Detection, Solely Responsible Agent Patterns, and Sandboxing}

A cooperative claim announces intended work. It prevents a collision only if
all conflicting effects respect an enforced admission rule. Different file
paths can refer to overlapping state, while disjoint edits in one file need
not conflict; the resource model must specify what overlaps.

One design gives a designated owner authority over a scarce resource while
other workers submit proposals. Another admits transactions after version
validation. Either still needs the effect boundary from the counter example.
Sandboxing limits reachable resources; it does not by itself choose the right
owner or make a stale update current. Chapters 1--3 develop these obligations
separately, so a claim record is never silently promoted into confinement.


\begin{figure}[tp]
  \centering
  % Reader question: which unsafe events can a supervisor disable?
% Claim: controllable events can be disabled; uncontrollable continuations remain.
% Evidence: upstream language containment and event-class automaton.
% Distinguish: safe target K, achievable controlled language and observation limits.
% Grammar: nested languages and named transitions, not a universal safety gate.
% Acceptance: state the controllability and observation conditions on K.
\noindent\begin{minipage}{11.1cm}\centering
\pdfigurehue{pdcobalt}%
\begin{tikzpicture}[pd figure,x=1cm,y=1cm]
  % Nested languages and a separate controllable-event boundary.
  \draw[pd neutral fill] (.10,1.00) rectangle (4.80,5.20);
  \node[pd label,anchor=north west] at (.25,5.05) {$L(G)$};
  \draw[pd focus fill] (1.05,1.60) rectangle (3.85,3.95);
  \node[pd label,anchor=north] at (2.45,3.78) {safe target $K$};
  \node[pd focus datum] at (2.35,2.55) {};
  \node[pd badge] at (2.70,2.90) {1};
  \node[pd breach datum] at (4.25,2.20) {};
  \node[pd badge] at (4.25,2.65) {2};
  \node[pd label] at (2.45,.55) {$L(S/G)\subseteq K\subseteq L(G)$};
  \draw[pd trust boundary] (6.30,4.00) rectangle (10.80,5.85);
  \node[pd kind] at (8.55,5.55) {$S$ disables $\Sigma_c$};
  \node[pd automaton state] (X0) at (7.20,2.65) {$x_0$};
  \node[pd automaton state] (X1) at (7.20,4.70) {$x_1$};
  \node[pd automaton state,draw=pderror,dash pattern=on 3.2pt off 2.2pt] (X3) at (9.80,4.70) {$x_3$};
  \node[pd automaton state] (X2) at (7.20,.65) {$x_2$};
  \draw[pd initial] (5.80,2.00) -- (X0);
  \draw[pd arrow] (X0) -- (X1);
  \node[pd badge] at (6.65,3.65) {1};
  \draw[pd breach arrow] (X1) -- (X3);
  \node[pd badge] at (8.55,4.15) {2};
  \draw[pd arrow] (X0) -- (X2);
  \node[pd label,anchor=west] at (7.90,1.65) {$\Sigma_u$};
\end{tikzpicture}
\par\smallskip
{\pdfiglabel\renewcommand{\arraystretch}{1.15}%
\begin{tabularx}{\linewidth}{@{}c p{2.0cm}p{1.0cm}>{\raggedright\arraybackslash}X@{}}
\toprule
 & \textbf{Event} & \textbf{Class} & \textbf{Supervisor}\\
\midrule
1 & file write & $\Sigma_c$ & Admitted: the trace stays inside $K$.\\
2 & process spawn & $\Sigma_c$ & Disabled: would leave $K$.\\
-- & tokens, signals, exit & $\Sigma_u$ & Cannot disable this continuation.\\
\bottomrule
\end{tabularx}}
\par\smallskip
{\pdfiglabelitalic\raggedright Containment requires implementability: permitted prefixes
must admit their uncontrollable continuations. Partial observation adds a
condition; an arbitrary safe target need not satisfy $L(S/G)=K$.\par}
\end{minipage}
%
  \pdwidecaption{\textbf{Control has a boundary.} A supervisor may disable a controllable write, but cannot simply disable a crash. A safe language must accommodate its permitted uncontrollable extensions; otherwise the supervisor must restrict earlier choices or report infeasibility.}[fig:wonham-supervisory]
\end{figure}

\subsection{Zombie Detection, Heartbeats, and State Salvage}

\begin{pdexample}{The expired worker resumes}
A worker holds fencing epoch 8 and pauses. After the declared lease deadline,
the authority grants a successor epoch 9. The old worker then resumes and sends
a write bearing epoch 8. The actuator must reject it. Reassigning a database
row alone would leave both workers able to write.
\end{pdexample}
A heartbeat is evidence of recent communication, not a proof of continued
health; a missed heartbeat can indicate delay rather than death. Lease expiry
is a policy decision. Fencing makes that decision effective at the resource.
Chapter 1 gives the state transitions, and Chapter 2 distinguishes epochs
from credential expiry.

Recovery must also preserve ownership of unfinished work. A successor can
inspect retained diffs, test records and obligations under a declared recovery
policy. It must not treat every untracked file as disposable or mistake a
missing response for an uncommitted operation. This reconnects the timeout
example to the work unit with which the chapter began.


\subsection{Closed-Loop Attestation and Execution Verification}

Figure~\ref{fig:pre-agent-loop} illustrates the closed agentic loop, stratified into three architectural tiers: the probabilistic Deliberation Layer, the confined Execution Layer, and the deterministic Governance Layer.

\begin{figure}[tp]
  \centering
  % Reader question: which responsibilities separate proposal, effect and evidence?
% Claim: governance checks admission; execution confines effects; receipts report.
% Evidence: three architectural tiers and two directed information flows.
% Distinguish: probabilistic proposal, mediated authority and observed outcome.
% Grammar: upstream three-band stack; architectural tiers, not temporal phases.
% Acceptance: locate both pre-effect admission and returned evidence.
\noindent\begin{minipage}{11.1cm}\centering
\pdfigurehue{pdcobalt}%
\begin{tikzpicture}[pd figure,x=1cm,y=1cm]
  \def\w{8.20}\def\h{1.40}
  \begin{scope}[xshift=1.35cm]
  \foreach \i/\name/\detail/\s in {
    2/deliberation/{LLM reasoning, context buffer, output validation}/{pd neutral fill},
    1/execution/{sandbox, host subprocess, spend cap}/{pd neutral fill},
    0/governance/{WAL ledger, capability checks, Merkle attestation}/{pd focus fill}}{
    \draw[\s] (0,{\i*\h}) rectangle (\w,{\i*\h+\h});
    \node[pd title,anchor=west] at (.20,{\i*\h+\h/2}) {\name};
    \node[pd label,anchor=east,text width=4.6cm,align=right] at ({\w-.20},{\i*\h+\h/2}) {\detail};}
  \draw[pd focus arrow] (-.75,{3*\h-.10}) -- (-.75,.10);
  \node[pd badge] at (-.75,{1.5*\h}) {1};
  \node[pd focus label,rotate=90] at (-1.15,{1.5*\h}) {intent};
  \draw[pd focus arrow] ({\w+.60},.10) -- ({\w+.60},{3*\h-.10});
  \node[pd badge] at ({\w+.60},{1.5*\h}) {2};
  \node[pd focus label,rotate=90] at ({\w+1.00},{1.5*\h}) {receipts};
  \end{scope}
\end{tikzpicture}
\par\smallskip
{\pdfiglabel\renewcommand{\arraystretch}{1.15}%
\begin{tabularx}{\linewidth}{@{}c p{1.2cm}p{2.65cm}>{\raggedright\arraybackslash}X@{}}
\toprule
 & \textbf{Flow} & \textbf{Direction} & \textbf{Carries}\\
\midrule
1 & intent & deliberation $\to$ governance & Proposed action, admitted before its mediated effect.\\
2 & receipts & governance $\to$ deliberation & Observed outcome and retained-evidence references.\\
\bottomrule
\end{tabularx}}
\par\smallskip
{\pdfiglabelitalic\raggedright Governance admits or refuses before execution.
Confinement can also block effects. A receipt records an observed outcome;
it does not establish semantic correctness.\par}
\end{minipage}
%
  \pdwidecaption{\textbf{Proposal, admission, execution and evidence.} The model proposes an action. The host checks and executes admitted effects, then records observations for the next decision. Correctness depends on complete mediation and the stated recovery protocol.}[fig:pre-agent-loop]
\end{figure}

\vspace{10pt}
\noindent\begin{minipage}{\linewidth}
  \centering
  \pdswissfigplate{fig-ch0-closed-loop}{0.88\linewidth}%
  \par\vspace{8pt}%
  {\color{pdcobalt}\rule{\linewidth}{1.2pt}}\par\vspace{5pt}%
  \raggedright
  {\color{pdswissred}\rule{4.5pt}{4.5pt}}\enspace%
  {\pdgrotesk\bfseries\scshape\fontsize{9}{11}\selectfont\color{hhink}%
    Architectural Plate 0.F\enspace\textperiodcentered\enspace Closed-Loop Attestation \& Governance}\par\vspace{4pt}%
  {\pdgrotesk\itshape\fontsize{9.5}{13}\selectfont\color{hhink}%
    ``Stratifying probabilistic deliberation from deterministic kernel attestation provides the authoritative bedrock for autonomous software engineering.''}\par\vspace{6pt}%
  {\color{hhgray!30}\rule{\linewidth}{0.4pt}}
\end{minipage}
\vspace{10pt}


A possible harness follows this control sequence. Its stored histories may be unbounded, so the sequence alone does not establish a finite-state model:
\begin{enumerate}[leftmargin=1.8em,itemsep=2pt]
  \item \textbf{Pre-Invocation Assembly:} The harness compiles system instructions, tool schemas, retrieved repository symbols, and recent conversation turns into the token buffer.
  \item \textbf{Inference:} The prompt is transmitted to the inference engine. Tokens stream autoregressively until a function-call termination token is detected.
  \item \textbf{Validation and Policy Gate:} The parser extracts tool arguments. The supervisor validates types against JSON-Schema definitions and verifies authority against local capability policies.
  \item \textbf{Sandboxed Execution:} The command executes inside an isolated runtime container (macOS Seatbelt, Linux Bubblewrap) under strict egress caps.
  \item \textbf{Kernel Witnessing:} Mutations, locks, and file reservations are submitted to the local kernel, appending a receipt to the audit log.
  \item \textbf{Observation Injection:} Raw output is sanitized, truncated to preserve token budget, and appended to the context buffer for the next deliberation step.
\end{enumerate}

\section{The Kernel Gap: Effects and Recovery}\label{pre:sec:kernel-gap}

Consider an agent with permission to propose any tool call. A useful guarantee must survive both an adverse proposal and a process crash. The following elementary result locates those proof obligations without assuming that stochastic generation always fails.

\begin{pdclaim}{Theorem}{Mediation and durable recovery boundaries}\label{thm:kernel-gap}
Let $I$ be a state invariant, $s_0\in I$, and $f(s,a)$ the effect of admitted action $a$.
\begin{enumerate}[leftmargin=1.5em,itemsep=1.5pt]
  \item \textbf{Complete mediation:} Suppose every state transition is checked at an indivisible admission-and-effect boundary against its current state, and every admitted pair satisfies $f(s,a)\in I$ whenever $s\in I$. Then every execution remains in $I$, regardless of the proposal policy. Conversely, if a reachable safe state has an unmediated admissible action leading outside $I$, universal safety over the admissible executions fails.
  \item \textbf{Conflicting updates:} If two actors may each implement an increment by a separate read and unconditional write, with no validation of the read version, their allowed interleavings include a lost update. Serializing the writes alone does not exclude that trace.
  \item \textbf{Crash recovery:} Suppose two pre-crash histories require different recovered obligation states, but erasure leaves identical surviving state and observations. No recovery procedure using only that information can recover the correct obligation state with certainty in both cases.
\end{enumerate}
\end{pdclaim}
\noindent\textit{Proof idea.} Preserve safety one transition at a time; exhibit a lost-update schedule; then compare two histories that recovery cannot distinguish.

\noindent\textit{Proof.} For (1), induction on the admitted transitions preserves $I$; the converse follows by appending the specified unsafe transition to its reachable prefix. For (2), both actors read 0 and then each writes 1, whereas two serial increments give 2. For (3), the recovery procedure receives identical input in the two cases, so its output distribution is identical; it cannot concentrate on two different required states.\hfill$\square$

Complete mediation may be implemented by a local gate, an authenticated remote permit checked before the effect, a verified controller, or another mechanism satisfying (1). A single-writer kernel is the construction developed here, not a unique mathematical solution. Disjoint or commuting operations need not share a mutex. Retained receipts or other distinguishing state are necessary for the recovery task in (3); refusing all further work can preserve some safety properties but does not recover the lost obligations. These distinctions are the basis for the chapters that follow.

\section{Chapter Map and Architectural Codas}\label{sec:prereq-map}

\subsection{Structure of the Volume}

The remainder of this textbook is organized into eight cumulative chapters across four distinct operational levels:

\begin{itemize}[leftmargin=1.8em,itemsep=1.5pt,topsep=1.5pt,parsep=0pt]
  \item \textbf{Chapter 1: The Single-Writer Kernel (Ground Truth).} We construct the minimal deterministic local kernel using SQLite WAL mode, establishing serializability, synchronous admission guards, and crash durability across process failures.
  \item \textbf{Chapter 2: The Anchor Protocol (Capability Delegation).} We design cryptographically verifiable capability tokens (Harbor Cards) and prove that delegated authority can only narrow at every hop.
  \item \textbf{Chapter 3: The Sealed Harbor (Confinement \& Leakage).} We construct isolated execution harbors that bound information leakage between distrustful principals under strict egress quotas.
  \item \textbf{Chapter 4: The Legible Swarm (Human Supervision).} We derive the information-theoretic cost of supervisory attention, proving how digests must retain evidence paths without drowning the operator in raw telemetry.
  \item \textbf{Chapter 5: From Spawn to Person (Durable Agency).} We define persistent agent identity, append-only outcome ledgers, and reputation mechanisms that survive process restarts.
  \item \textbf{Chapter 6: The Harbor Economy (Markets \& Clearing).} We establish bilateral clearing, underwriting reserves, and dynamic pricing for scarce host resources.
  \item \textbf{Chapter 7: The Bonded Commons (Game-Theoretic Governance).} We analyze claim signaling under stated monitoring and payoff assumptions, and separate deterrence from accounting and adjudication.
  \item \textbf{Chapter 8: The Federated Harbor (Cross-Machine Evidence).} We study witness logs and topology-dependent consistency checks while preserving local admission. A vanishing residual does not certify every fact or a globally consistent swarm.
\end{itemize}

\section{Review of the key ideas}
The opening counter separated a proposed update from an admitted transition.
The failed patch separated an unsuccessful attempt from an abandoned obligation.
The timeout separated a missing response from a missing effect. These are the
same discipline applied to actions, goals and observations: name the states
that matter, retain evidence that distinguishes them, and place a check where
it can still affect the outcome.

Commitment theory helps describe what persists; message kinds distinguish a
request from an acceptance; incentives depend on detection and enforceable
consequences. None replaces the others. The memory and attention calculations
explain why a prompt cannot carry every record. Chapters 1--3 establish
transaction, authority and disclosure boundaries. Chapters 4--8 use those
boundaries for supervision, succession, exchange and federation.

\section{Exercises}
\pdexercisesfor{Foundations}{Proposals, persistence and evidence}
\begin{pdexercise}[kind=Check,rating=1]{ex:prereq-increment}
Two workers each read counter 4 and propose writing 5. What final value follows
if only their writes are serialized? What value would two serial increments
give? State the missing check.
\end{pdexercise}
\begin{pdsolution}{ex:prereq-increment}
The writes give 5; serial increments give 6. Either include each read and update
in one serial transaction, or validate the read version atomically with the
write and retry the rejected operation from a fresh read.
\end{pdsolution}

\begin{pdexercise}[kind=Check,rating=1]{ex:prereq-risk}
A deviation gains 10 only when undetected; detection loses the gain and incurs
sanction 100. Find the smallest detection probability making its expected
incremental payoff nonpositive. Why is audit selection alone insufficient?
\end{pdexercise}
\begin{pdsolution}{ex:prereq-risk}
$10(1-p)-100p\le0$ gives $p\ge1/11$. If detection happens only during an audit, $p$ equals audit-selection
probability times successful detection conditional on selection. The sanction
must actually be enforceable.
\end{pdsolution}

\begin{pdexercise}[kind=Trace,rating=2]{ex:prereq-timeout}
A test-payment request times out. Give two histories consistent with the same
observation. Explain what an operation identifier and a durable receipt can
resolve, and what they cannot establish about the quality of the tested code.
\end{pdexercise}
\begin{pdsolution}{ex:prereq-timeout}
The request may never have arrived, or it may have committed before the response
was lost. A positive receipt identifies the committed operation. An absent receipt
alone does not prove that the request never arrived: it may still be in flight.
Safe retry therefore needs atomic duplicate suppression at admission, or a
terminal rejection/cancellation contract excluding a later first commit.
The receipt establishes the recorded operation, not arbitrary semantic
correctness of the tested code.
\end{pdsolution}

\begin{pdexercise}[kind=Trace,rating=2]{ex:prereq-knowledge}
A hidden bit is uniform on $\{0,1\}$. Both cases emit ``ready,'' but an exit code
then equals the bit. List the remaining possible states before and after
observing exit code 0. Does hiding the file alone preserve secrecy?
\end{pdexercise}
\begin{pdsolution}{ex:prereq-knowledge}
Before the code, both states remain possible; afterward only bit 0 remains.
The exit code reveals the bit even though the file is unreadable. A confinement
argument must include that channel.
\end{pdsolution}

\begin{pdexercise}[kind=Open,rating=3]{ex:prereq-recovery-contract}
Design a recovery contract for Alice's repair after a worker pauses with an
unreported external effect. Specify authoritative state, an ambiguous pair of
histories, the distinguishing evidence, a stale-worker refusal rule and one
failure your contract cannot repair. Explain how a replacement worker learns
what remains owed without repeating a settled effect.
\end{pdexercise}
\begin{pdsolution}{ex:prereq-recovery-contract}
A defensible design identifies the effect's actuator, operation identifier,
commit and receipt boundary, recovery lookup, successor epoch and obligation
record. Exhibit both a retry and a stale-worker trace. If an external service
cannot deduplicate or report the operation, acknowledge that exact recovery
may be impossible and specify escalation instead of inventing certainty.
There is no unique protocol or closed-form solution.
\end{pdsolution}

\section{History and references}
The chapter draws its distinctions from several traditions: Bratman and
Rao--Georgeff for intentions, actor and CSP models for communication,
communicative-act theory for message intent, and supervisory control for
what a gate can enforce. Arrow and mechanism design identify limits on
collective decisions; Ostrom supplies institutional questions rather than a
software correctness theorem. The cited planning and learning systems concern
proposal generation and evaluation. Their benchmarks are not certificates
for the host mechanisms developed next.


\enlargethispage{2.5\baselineskip}
\vspace{-2pt}
\pdchapterhandoffprereq

\clearpage
